\documentclass{xdupgthesis}

% XDUTS核心配置
\xdusetup{
  style = {
    customize-los = false,
    customize-loa = false,
  },
  info = {
    graduate-type         = 硕士,
    degree-type           = 专业,
    % 核心必填信息
    title                 = {深度学习多模态眼动事件检测},
    title*                = {Multi-Modal Eye Movement Event Detection Based on Deep Learning},
    author                = {刘畅},
    author*               = {Zhiyong Yu},
    department            = {电子工程学院},
    major                 = {电子信息},
    major*                = {Electronic Information},
    sub-major             = {智能感知与信号处理},
    student-id            = {23021211914},
    supervisor            = {郑洋},
    supervisor*           = {Yang Zheng},
    supv-title            = {讲师},
    supv-title*           = {Lecturer},
    supervisor-enterprise = {},
    supervisor-enterprise*= {},
    supv-ent              = {},
    supv-ent*             = {},
    supv-ent-title        = {},
    supv-ent-title*       = {},
    supervisor-title      = {讲师},
    supervisor-title*     = {Lecturer},
    supervisor-enterprise-title  = {},
    supervisor-enterprise-title* = {},
    % 学科与保密信息
    domain                = {新一代电子信息技术（含量子技术等）},
    clc                   = {TP391.41},
    secret-level          = {公开},
    submit-date           = {2026-03-01},
    % 关键词
    keywords              = {眼动事件检测, 深度学习, 多模态融合, Mamba网络, 时频分析},
    keywords*             = {Eye Movement Event Detection, Deep Learning, Multi-Modal Fusion, Mamba Network, Time-Frequency Analysis},
    % 附件与声明文件
    abstract              = {chapters/abstract-zh.tex},
    abstract*             = {chapters/abstract-en.tex},
    acknowledgements      = {chapters/acknowledgements.tex},
    bio                   = {chapters/biography.tex},
    appendix              = {chapters/appendix1.tex},
    statement-scan        = {},
    statement-sign        = {},
    los                   = {chapters/los.tex},
    loa                   = {chapters/loa.tex},
    % 参考文献
    bib-resource          = {references.bib},
  }
}

% -------------------------
% 常用包（按功能分组）
% -------------------------
\usepackage{graphicx}
\graphicspath{{images/}}

\usepackage{amsmath}
\usepackage{amsfonts}
\usepackage{amssymb}

\usepackage{booktabs}
\usepackage{multirow}
\usepackage{makecell}
\usepackage{tabularx}
\usepackage{threeparttable}

\usepackage{subcaption}

\usepackage{xcolor}
\usepackage{pifont}

\usepackage{enumitem}

% 浮动体控制（你已在用：建议保留）
\usepackage[section]{placeins} % 提供 \FloatBarrier，并阻止浮动体跨 section
\usepackage{flafter}          % 防止浮动体出现在源码位置之前

% 超链接与智能引用：hyperref 必须在 cleveref 前面
\usepackage{hyperref}
\usepackage{cleveref}

% -------------------------
% 全局宏：模型名（你原本就有）
% -------------------------
\newcommand{\ModelTF}{GazeFusion}
\newcommand{\ModelMM}{GazeGuide}

% -------------------------
% 全局宏：地磁章节统一口径（按你第2章需要）
% -------------------------
\newcommand{\Fs}{\ensuremath{50\,\mathrm{Hz}}}
\newcommand{\MagUnit}{\ensuremath{\mathrm{nT}}}
\newcommand{\SmoothLen}{\ensuremath{7}}        % L_s
\newcommand{\BgWinPts}{\ensuremath{500}}       % N_0 (10 s @ 50 Hz)
\newcommand{\ProbFloor}{\ensuremath{10^{-6}}}  % p_min

% 兜底：即使 placeins 被注释掉，也不至于 \FloatBarrier 报错
\providecommand{\FloatBarrier}{}

% -------------------------
% 文档开始
% -------------------------
\begin{document}


\chapter{绪论}

\section{研究背景与意义}

近年来，随着眼动追踪技术的快速发展，眼动分析已成为神经科学、认知心理学、临床医学、教育学及人机交互等领域的重要研究手段，在神经退行性疾病诊断、情绪认知评估、人机交互设计、教育系统智能化以及广告效应评估等方面具有广泛应用前景。例如，阿尔茨海默病患者在反跳视任务中表现出的眼跳潜伏期延长和错误率升高，可作为疾病早期诊断的潜在生物标志物 \cite{readman2021potential}；抑郁症患者则存在对负性刺激的凝视偏好和对正性刺激的回避行为，反映其情绪加工过程中的认知偏差 \cite{noyes2023eye}。在智能驾驶领域，眼动分析可用于评估驾驶员注意力状态，提高行车安全；在虚拟现实技术中，眼动追踪为更自然的交互方式提供了可能。

眼动事件检测是眼动分析的核心环节，其目标是将原始眼动数据分割为注视（Fixation）、眼跳（Saccade）、扫视后震荡（PSO）和平滑追踪（Smooth Pursuit）等具有明确生理意义的事件类型。准确的事件检测是提取有效眼动特征、保证实验结论可靠性的前提，直接关系到眼动分析的后续应用效果。然而，眼动信号具有非平稳性、低信噪比和事件边界模糊等特点，传统检测方法基于手动设定的特征和阈值，通过分析眼动速度、位置离散度等定量指标进行事件分类，过度依赖人工经验和数据分布假设，难以适应个体差异和动态的环境变化，且不同场景下的检测精度和稳定性不尽理想。

随着机器学习和深度学习方法的发展，学术界逐渐尝试通过自动特征提取与模式学习的方式提升眼动事件检测的准确性和鲁棒性。深度学习技术，特别是基于卷积神经网络（CNN）、长短时记忆网络（LSTM）以及更具时序捕捉能力的Transformer架构，使得眼动事件检测进入了端到端自动化分类的新时代。然而，现有的深度学习方法仍面临诸多挑战：在处理长距离依赖的时序特征和捕捉多模态动态信息方面能力有限，存在全局信息捕捉与局部边界定位的矛盾，且对眼动序列的空间结构特征和多源模态信息（如瞳孔变化、场景图像）的融合利用仍不充分，导致其在复杂情境下的表现力受限。

本研究选题的重要性不仅在于推动眼动事件检测技术的前沿发展，还在于探索更具鲁棒性与适用性的多模态融合方法，为眼动事件检测开辟新路径。多模态数据融合的引入为解决单一模态在捕捉复杂时空信息上的局限性提供了新思路，视频模态的数据能够补充序列数据在动态环境下的细节信息，从而提高对复杂场景中眼动事件的捕捉能力。此外，为了应对眼动序列中长距离时序依赖难以捕捉的问题，结合Transformer及其他时序特征增强模块以提高模型的全局感知能力，对于提升眼动事件检测的准确性和模型的泛化能力具有重要意义。

本研究不仅具有重要的理论价值，还对人类视觉认知的实际应用具有深远意义。通过提升眼动事件检测的精确度和环境适应性，可以进一步推动在教育培训、智能驾驶、虚拟现实等多领域的创新应用，为复杂交互环境下的人机系统设计提供技术支撑，同时也为深入理解人类视觉认知过程提供新的研究视角和方法。

\section{国内外研究现状}

眼动事件检测作为视觉认知、注意力研究及行为分析的重要方法，一直是国内外研究的热点领域。随着数据采集和计算技术的发展，研究方法从传统的特征工程逐渐过渡到基于机器学习和深度学习的自动化模式识别方法。以下将从传统方法、机器学习方法和深度学习方法三个层次，梳理国内外在眼动事件检测方面的研究进展及其面临的挑战。

\subsection{传统阈值法}
传统的眼动事件检测方法主要依赖于手动定义的特征和阈值，通过对眼动特性指标（如速度、加速度和位置离散度）的解析，实现对注视、扫视等事件的分类。这类方法中，经典的基于离散度（I-DT）和速度（I-VT）的方法较为常用\cite{salvucci2000identifying}。

I-VT方法通过计算点对点速度来识别事件，然后根据该速度的值将事件分类为注视或扫视。经典的I-VT方法旨在将所有眼动追踪输入数据仅分类为注视点和眼跳。其他事件类型，如平滑追逐、后扫视振荡和噪声，则不予考虑。I-DT则根据后续样本坐标的色散或扩散距离对注视点和扫视点进行分类。当样本位于空间有限的区域内时，将凝视数据识别为属于注视点，否则全部识别成眼跳。

近年来，研究人员进一步发展了IVVT（结合速度和离散度阈值）和IVMP（基于速度和运动模式）等改进算法，增加了对平滑追踪事件的识别能力。尽管传统方法简单易行，在采样频率较低的情境下表现良好，但其在处理动态和复杂场景中表现不够理想。由于手动阈值的设定往往依赖于预先定义的经验规则，这类方法在个体间具有较大的不稳定性。此外，阈值设定缺乏标准，往往难以在不同实验设置间保持检测的一致性。随着眼动数据采集分辨率的提升，传统方法难以捕捉到高维特征和长距离依赖关系，逐渐暴露出局限性。

\subsection{机器学习方法}
通过训练分类器（如随机森林、隐马尔可夫模型）自动学习特征与事件类型的映射关系 \cite{zemblys2017eye,pekkanen2017new}。随着计算能力和数据处理技术的进步，研究者们开始尝试将机器学习应用于眼动事件检测，以期提升模型的鲁棒性和自动化水平。

机器学习方法通常通过提取与眼动事件相关的高维特征并引入分类模型，实现事件分类。例如，Vidal等人提出了一种基于眼动形状特征的平滑追踪检测方法，该方法利用特征提取与机器学习分类模型结合的方式，能够自动化地识别平滑追踪事件，但无法分类其他类型的眼动事件。

在多类眼动事件分类方面，基于分段线性回归的Naive Segmented Linear Regression (NSLR) 方法展示了较强的稳定性和适用性。该方法结合隐马尔可夫模型（HMM），通过分段拟合不同眼动事件的特征，能够在噪声环境中保持较高的分类精度。研究显示，NSLR能够有效识别注视、扫视、平滑追踪和后扫视振荡（PSO），具有较高的事件区分度。此外，Pekkanen与Lappi等人利用随机森林算法构建了眼动事件分类模型，并证明了其在不同采样频率和噪声水平下的鲁棒性。

尽管机器学习方法在眼动事件检测上取得了显著进展，但其依赖于复杂的特征提取和前期数据处理，且通常需要大量带标签的数据集来训练模型，这在实际应用中增加了实现的难度。此外，传统机器学习模型难以捕捉长时序依赖和眼动事件间的动态特征，表现出一定的局限性。

\subsection{深度学习方法}
基于神经网络自动提取特征，实现端到端事件检测，主要包括：

深度学习技术凭借其强大的自动特征学习和模式识别能力，在眼动事件检测中展现出卓越的应用潜力。相较于传统方法和机器学习方法，深度学习能够在端到端的框架中自动提取高维特征，且无需大量人工设定的规则和参数，极大地提升了模型的适用性与推广性。

1. 卷积神经网络（CNN）：擅长捕捉局部时空特征，但 receptive field 有限，难以建模长时依赖 \cite{hoppe2016end,birawo2022review}。在深度学习方法中，CNN被广泛应用于眼动序列的特征提取。Hoppe与Bulling提出了基于CNN的眼动事件检测方法FFT-CNN，该方法通过频域特征变换解决了传统方法中手动阈值设定的问题，实现了眼动事件的自动检测。

2. 循环神经网络（RNN/LSTM）：适合序列建模，但存在梯度消失问题，长序列处理效率低 \cite{goltz2019exploring,li2021attention}。然而，CNN缺乏对时间依赖的记忆能力，难以捕捉序列数据的时序特征，为此研究者逐渐引入了LSTM、Bi-LSTM等具备长时序记忆功能的网络。例如，gazeNet模型结合了CNN和Bi-LSTM层，在眼动数据的时间特征提取方面取得了良好的效果，并提出了一种生成数据的方式以增强模型的泛化能力。

3. 混合模型（CNN+RNN/Transformer）：融合局部特征提取与全局依赖建模能力 \cite{zemblys2019gazenet,suhari2023improving}，但部分模型存在计算复杂度高或边界定位精度不足的问题。为了提升模型对序列长距离依赖的处理能力，研究者开始将Transformer应用于眼动事件检测任务。Transformer模型通过自注意力机制在全局范围内建立特征关联，使得模型能够有效捕捉序列中的长距离依赖关系。近年来，基于Transformer结构的网络逐渐发展为眼动事件检测中的新兴技术。Wang等人提出的Two-Stream Vision Swin Transformer模型，通过融合光流和视频帧信息，实现了高精度的眼动事件检测。该方法在动态视觉场景下展示了较强的检测能力，为解决头动和光照变化等复杂问题提供了有效途径。

4. 多模态融合方法：初步尝试结合时域、频域特征，但对空间结构特征和多源模态信息（如瞳孔变化、场景图像）的融合利用仍不充分。

深度学习方法虽然展现出极大的优势，但当前研究中仍存在若干挑战。首先，CNN和LSTM等模型在处理长序列数据时感受野受限，难以捕捉全局的时序关系。其次，深度学习模型的计算复杂度高，模型参数众多，对硬件资源的要求较高。此外，由于眼动事件数据的标注成本高昂，数据集匮乏仍然是深度学习方法面临的主要问题之一。为应对这一问题，研究者们逐渐探索数据增强和合成数据生成技术，以扩充数据量，提高模型泛化能力。

综上所述，眼动事件检测研究的发展历程显示出传统方法、机器学习方法和深度学习方法的演变。传统方法在低维眼动特征的阈值分类上取得了一定成效，但在复杂情境下表现受限；机器学习方法引入了自动化分类模型，提升了鲁棒性，但仍受限于手动特征工程的需求。深度学习方法则通过自动特征提取与端到端训练模式，展现出卓越的检测能力。然而，现有深度学习方法在长距离依赖捕捉、动态环境适应及数据需求等方面仍有提升空间。

\section{研究目标与内容}

\subsection{研究目标}
本研究旨在构建多模态融合的眼动事件检测方法，解决当前方法在动态复杂场景中的适应性不足问题，具体目标为：
1. 提出轻量高效的长距离依赖特征捕捉方法，提升时序建模能力；
2. 设计基于坐标嵌入的特征提取技术，增强空间位置关系理解；
3. 建立序列与视频数据融合策略，提高复杂情境下的事件检测精度。

通过实现以上目标，推动眼动行为检测技术进展，增强其在真实动态场景中的实用性和鲁棒性。

\subsection{研究内容}
本研究将围绕以下几个核心内容展开：

基于长距离依赖增强的时序特征建模方法：针对现有眼动事件检测模型在长时序数据处理中的局限性，拟提出一种结合Transformer与Mamba结构的时序特征增强模块，以捕捉全局时序依赖。通过设计轻量化的注意力机制，使得模型能够高效地处理长序列数据，提升分类准确率和鲁棒性。

坐标嵌入技术的优化与数据增强：为解决现有模型对空间关系理解不足的问题，提出一种基于坐标嵌入的特征提取方法，将眼动序列中的(x, y)坐标转化为高维特征向量。该嵌入方法将采用几何位置编码、卷积嵌入与图神经网络等多种方式，通过综合比对实验优化特征提取效果，以捕捉眼动事件中各坐标点的相对空间关系。

在数据增强与模型优化方面，探索不同的数据增强方式，包括时间扰动、镜像翻转、旋转等操作，以提升模型对不同噪声条件下的鲁棒性。同时，通过引入监督和自监督学习策略优化模型训练过程，提升模型对小样本数据的适应性。

多模态融合策略的设计与实现：针对单一序列数据在动态场景下表现不足的问题，引入视频模态进行多模态数据融合。通过结合序列数据与视频数据的特征，增强模型对动态事件的理解和场景适应能力。该部分将设计信息融合模块，通过双流网络、加权特征拼接、时空注意力等方法实现多模态信息的深度融合，确保模型在复杂场景中具备较高的泛化能力。

\subsection{拟解决的关键问题}

长距离依赖特征捕捉能力不足：现有模型在处理眼动序列数据中的长时序依赖时存在局限性，难以捕捉全局时序关系。为此，本研究拟通过Transformer等注意力机制增强时序建模能力，解决长序列数据处理中的有效信息提取问题。

空间关系特征的表达欠缺，数据缺乏与样本多样性不足：传统眼动事件检测方法中，(x, y)坐标点作为绝对序列数据输入模型，未能有效表征各坐标点间的相对关系。本研究提出的坐标嵌入技术将解决此问题，通过将眼动事件中的空间特征转化为高维特征表示，提升模型对空间关系的感知能力。并计划使用多种数据增强方法与合成数据生成技术，扩充数据量以提高模型的泛化能力。

动态场景中的模态信息不足：单一序列数据在复杂动态场景中的表现有限，难以适应环境光照变化、头部运动等干扰因素。为此，本研究拟通过引入视频模态，实现多模态信息融合，增强模型在动态场景下的适应性与精度。

\section{论文结构安排}

本论文共分为6章，各章节内容安排如下：

第1章 绪论。本章首先介绍眼动事件检测的研究背景与意义，阐述眼动追踪技术在神经科学、人机交互、心理学等领域的应用价值。接着系统梳理眼动事件检测的发展历程，从传统阈值法、机器学习方法到深度学习方法的演变过程，分析现有方法在长距离依赖捕捉、空间关系理解及动态场景适应等方面存在的局限性。最后明确本研究的目标、核心内容、拟解决的关键问题及技术路线，为后续研究工作奠定基础。

第2章 相关理论与方法。本章介绍眼动事件检测的基础理论与关键技术。首先详细阐述眼动事件的定义、分类及特征，包括注视、扫视、平滑追踪等基本眼动事件的生理机制与数据表现。其次，系统介绍深度学习核心技术，包括卷积神经网络（CNN）、循环神经网络（RNN/LSTM）、Transformer及Mamba等模型的原理与应用。最后，讲解时频分析方法与空间编码技术，包括短时傅里叶变换（STFT）、坐标嵌入、位置编码等，为后续模型设计提供理论支撑。

第3章 双流时频眼动事件检测。本章提出基于时频分析的双流眼动事件检测框架\ModelMM{}。首先，设计基于短时傅里叶变换的时频特征提取模块，将一维眼动序列转化为二维时频图像，以捕捉眼动信号的时频域信息。其次，构建双流网络结构，分别处理原始序列数据和时频图像数据，并通过注意力机制实现特征融合。最后，在公开数据集上进行实验验证，包括参数调优、消融实验及与现有方法的对比分析，验证该框架的有效性与优越性。


第4章 眼动序列空间编码研究。本章提出基于坐标嵌入的眼动序列空间编码方法。首先，分析传统眼动数据处理中空间信息丢失的问题，设计多种坐标嵌入方案，包括几何位置编码、卷积嵌入与图神经网络嵌入等。其次，构建空间-时序融合模型，将嵌入后的空间特征与时序特征进行深度融合，提升模型对眼动事件空间关系的理解能力。最后，通过对比实验验证不同嵌入方法的效果，并分析模型在不同噪声条件下的鲁棒性。

第5章 多模态眼动事件检测。本章构建多模态眼动事件检测系统，结合序列数据与视频模态数据提升动态场景适应性。首先，建立多模态数据集，同步采集眼动序列数据与场景视频数据，并进行数据预处理与标注。其次，设计双流多模态融合网络，分别提取序列特征与视频特征，并通过时空注意力机制实现信息融合。最后，在动态场景数据集上进行实验评估，分析多模态融合对复杂场景下事件检测精度的提升效果。

第6章 总结与展望。本章对全文研究工作进行总结，梳理本研究的主要贡献与创新点，包括长距离依赖特征捕捉方法、坐标嵌入技术及多模态融合策略等。同时，分析研究工作中存在的局限性，如数据样本多样性不足、模型实时性有待提升等。最后，展望眼动事件检测领域的未来研究方向，包括大模型应用、跨模态融合、边缘计算优化等，为后续研究提供参考。






% ============================================================
% Chapter 2: 地磁车辆事件检测基本原理与系统设计
% 说明：
% 1) 本章默认在导言区已定义以下全局宏：
%    \Fs \MagUnit \SmoothLen \BgWinPts \ProbFloor
% 2) 本章使用 \FloatBarrier（导言区已加载 placeins[section] 或提供兜底）
% ============================================================

\chapter{地磁车辆事件检测基本原理与系统设计}
\label{ch:ch2_event_detect}

\section{引言}\label{sec:ch2_intro}

随着智慧交通系统的发展，道路侧车辆信息的准确获取是交通状态监测、运行评估与管理优化的重要基础\cite{its_overview,traffic_monitoring}。
车辆通过事件、车型类别以及停车占用状态等数据可为车流统计、运行状态评估与交通组织优化提供支撑\cite{vehicle_event_application}。
地磁传感器具有易于路侧与地面部署、功耗低以及适合分布式长期运行等特点，可用于道路及停车场景的车辆信息连续采集与上报\cite{magnetic_sensor_traffic}。

然而，实际道路环境中的地磁观测除车辆扰动外，还会受到温度漂移、慢变背景漂移、电磁噪声与偶发尖峰等因素影响，导致基线随时间变化并出现异常波动。
同时，车辆速度变化会引起扰动波形在时间轴上的伸缩，使事件边界定位与后续特征提取的稳定性下降。
因此，在开展车型分类与停车占用判定之前，有必要建立鲁棒的车辆检测与事件分割底座，并明确端侧处理、通信上报与平台汇聚的链路流程，
以保证后续算法研究的可复现性与数据口径一致性。

本文采用单点三轴地磁传感器，以采样频率 $f_s = \Fs$ 获取连续磁场序列。
为便于表达微小扰动与差分变化，本文将磁场强度及其差分量统一以 \MagUnit\ 表示（$1\,\mu\mathrm{T}=1000\,\mathrm{nT}$）。

\paragraph{工程约束与设计取舍（采样率与上报策略）}
本系统面向道路侧长期在线部署，端侧节点通常由电池供电，并通过低功耗广域无线链路（如 LoRa）进行数据上报。
若将三轴原始波形持续回传，即使仅按 16 bit 整型量化，上报数据量也约为
$3 \times f_s \times 2 \approx 300\,\mathrm{B/s}$，
折算为约 $1.08\,\mathrm{MB/h}$、$25.9\,\mathrm{MB/day}$。
在低功耗链路下，这将显著增加空口占用、节点能耗与丢包风险。
因此，系统采用“端侧事件提取 + 事件级摘要上报”的分工：端侧实时输出事件边界与驻留时长等摘要字段，
云端负责汇聚、存储、可视化与统计分析。
原始三轴事件片段仅在实验采集与离线标注阶段获取，用于后续章节算法验证，从而避免常态运行时带宽与功耗负担。

本章首先在第~\ref{sec:ch2_principle}~节介绍地磁车辆检测的基本原理，并在第~\ref{sec:vehicle_detection}~节给出基于差分特征、概率融合与有限状态机的车辆检测与事件分割方法，
从连续数据流中稳定定位车辆进入与离开边界并形成标准化事件片段。
随后在第~\ref{sec:ch2_platform}~节说明车辆信息检测平台的组成与功能，形成端侧事件提取、链路传输、云端管理与前端可视化的闭环系统；
在第~\ref{sec:ch2_module}~节给出端侧检测模块的组成、选型依据与数据口径。
上述内容为后续车型三分类与稳定点停车占用判定提供统一的数据入口与系统支撑。

\section{地磁车辆检测基本原理}\label{sec:ch2_principle}

\subsection{车辆检测算法基本原理}\label{sec:ch2_basic}

地磁车辆检测的基本依据在于：在一定空间范围内，地球磁场的强度与方向在短时间尺度上具有相对稳定性，可近似视为道路场景的背景磁场。
当车辆进入传感器邻域时，由于车体中铁磁性材料及其磁化效应，会改变局部磁场分布，从而在传感器输出的三轴磁场序列中形成可观测的扰动特征。
因此，对地磁场进行连续采样并监测其扰动变化，可实现车辆通过事件的检测与事件片段分割。

对于三轴地磁传感器，在离散采样时刻 $k$ 的输出可记为
\begin{equation}
    \mathbf{B}(k) = [B_x(k),\,B_y(k),\,B_z(k)]^{T}.
\end{equation}
其中 $B_x(k),B_y(k),B_z(k)$ 的单位均为 \MagUnit。
本文约定传感器坐标系与安装方向保持一致，其中 $X$ 轴沿车辆行驶方向，$Y$ 轴沿车道横向，$Z$ 轴垂直向上（以实际安装标定为准）。

\paragraph{车辆等效磁偶极子解释与扰动衰减特性}
为从物理层面对“扰动具有空间局部性、且三轴响应差异显著”给出解释，可将车辆及其磁化效应等效为有限数量磁偶极子的叠加。
对单个磁偶极子，距传感器位置向量为 $\mathbf{r}$（$r=\|\mathbf{r}\|$），偶极矩为 $\mathbf{m}$，其静磁场可写为
\begin{equation}
    \mathbf{B}_{dip}(\mathbf{r})=\frac{\mu_0}{4\pi r^3}\Big(3(\mathbf{m}\cdot \hat{\mathbf{r}})\hat{\mathbf{r}}-\mathbf{m}\Big),
\end{equation}
其中 $\hat{\mathbf{r}}=\mathbf{r}/r$。
该表达式表明扰动幅值随距离呈 $r^{-3}$ 量级衰减，因此车辆扰动主要集中在传感器邻域；同时，偶极矩方向、车辆姿态与传感器坐标系的相对关系
会导致 $X/Y/Z$ 三轴分量呈现不同幅值与波形形态。
这也是本文采用“三轴联合特征 + 概率融合”而非单轴阈值判决的直接依据。

当车辆驶入检测区域时，其金属结构会使周围磁感线发生畸变；由于车辆相对传感器的位置、姿态与结构差异，各轴响应幅度与波形形态可能不同，
从而使三轴输出出现不同程度的偏移与波动。车辆经过传感器邻域时的磁场扰动示意见图~\ref{fig:mag_disturbance}。

\begin{figure}[!htbp]
    \centering
    \includegraphics[width=0.86\linewidth]{fig_mag_disturbance}
    \caption{车辆经过地磁传感器时的磁场扰动示意图}
    \label{fig:mag_disturbance}
\end{figure}

车辆检测算法的目标可表述为：在连续采样的三轴磁场序列中识别车辆进入与离开检测区域的边界，并输出单车事件片段
\begin{equation}
    [k_{\mathrm{in}},\,k_{\mathrm{out}}),
\end{equation}
其中 $k_{\mathrm{in}}, k_{\mathrm{out}} \in \mathbb{Z}$ 为采样点索引；
$k_{\mathrm{in}}$ 表示车辆扰动进入检测区域的起始索引，
$k_{\mathrm{out}}$ 表示扰动结束并恢复至无车状态后的第一个索引。
相应驻留点数与驻留时长分别为
\begin{equation}
    N_{\mathrm{occ}} = k_{\mathrm{out}} - k_{\mathrm{in}},\qquad
    T_{\mathrm{occ}} = \frac{N_{\mathrm{occ}}}{f_s}.
\end{equation}
若需表示最后一个有车采样点，可记为 $k_{\mathrm{last}} = k_{\mathrm{out}} - 1$；
对应物理时间可由 $t = k/f_s$ 换算得到。
事件片段是后续车流统计、车型识别与停车占用判定等上层任务的基础输入。

为便于分析与算法设计，通常将观测序列表示为背景磁场、车辆扰动与噪声干扰的叠加
\begin{equation}
    \mathbf{B}(k) = \mathbf{B}_{bg}(k) + \mathbf{B}_{veh}(k) + \mathbf{w}(k),
\end{equation}
其中 $\mathbf{B}_{bg}(k)$ 为环境背景分量，$\mathbf{B}_{veh}(k)$ 为车辆引起的扰动分量，$\mathbf{w}(k)$ 为测量噪声及偶发电磁干扰等随机项。
车辆检测的关键在于构造能够突出 $\mathbf{B}_{veh}(k)$、同时尽量抑制 $\mathbf{B}_{bg}(k)$ 与 $\mathbf{w}(k)$ 影响的特征量，
并在该特征量上设计稳定的判决机制，实现事件边界的可靠定位。

在实际道路环境中，背景磁场并非严格常数：传感器温漂、零偏变化以及外界电磁环境变化会导致 $\mathbf{B}_{bg}(k)$ 随时间发生慢变漂移，
从而造成基线水平变化。长期运行条件下的三轴磁场慢变漂移示例见图~\ref{fig:mag_drift}。
若算法依赖固定阈值或依赖稳定基线估计，则基线偏移可能引起误触发或漏检；
同时在车流密集场景下，用于更新基线的稳定区间可能难以获得，进一步影响检测精度。
因此，面向长期运行的车辆检测算法需具备对慢变漂移的鲁棒性，并在存在瞬时尖峰与离群点等异常扰动时保持稳定判决。

\begin{figure}[!htbp]
    \centering
    \includegraphics[width=0.86\linewidth]{fig_mag_drift}
    \caption{长时间运行条件下三轴磁场序列的慢变漂移示例（纵轴单位为 \MagUnit）}
    \label{fig:mag_drift}
\end{figure}

此外，车辆速度变化会导致车辆扰动波形在时间轴上呈现拉伸与压缩现象，即扰动形态在空间上具有相似性，但在时间尺度上持续长度不同。
该特性会降低固定时间窗或固定持续时间阈值策略在不同速度条件下的适用性，并可能对后续特征提取与识别造成不利影响。
因此，检测阶段应通过鲁棒特征构造与时序判决机制降低速度变化与噪声的影响；在完成事件分割后，可进一步通过对齐与规范化等策略提升后续算法对速度变化的适应性。

综合上述分析，本文检测底座采用三轴共同参与的处理路径：以三轴磁场序列为输入，
通过构造鲁棒特征（如三轴一阶差分及其平滑）突出车辆扰动并抑制慢变背景，
并在融合输出上引入有限状态机与连续点数确认机制抑制瞬时尖峰与短时回落，从而稳定输出车辆事件片段 $[k_{\mathrm{in}}, k_{\mathrm{out}})$。
综合特征量的一般形式可写为
\begin{equation}
    F(k) = G\big(B_x(k), B_y(k), B_z(k)\big),
\end{equation}
其中 $G(\cdot)$ 表示特征构造算子。
在第~\ref{sec:vehicle_detection}~节中，将进一步给出基于一阶差分的特征构造方式以及状态机判决流程，以在复杂道路环境下稳定定位 $k_{\mathrm{in}}$ 与 $k_{\mathrm{out}}$。

\subsection{基于一阶差分与状态机的车辆检测算法}\label{sec:vehicle_detection}

本节给出本文车辆检测与事件分割算法。地磁传感器以采样率 $f_s = \Fs$ 连续采集三轴磁场序列，
车辆经过传感器邻域时会引起局部磁场畸变，使输出产生显著扰动。
在实际道路环境中，环境磁场慢变漂移、传感器温漂以及偶发电磁干扰会导致信号基线随时间变化；
在 \Fs\ 条件下，车辆扰动波形在时域上更易出现时间尺度拉伸与压缩，并可能因车底弱磁区域出现短时回落。
若直接在原始磁场幅值上进行阈值判决，容易产生误检与漏检。
为此，本文构建由三轴一阶差分特征提取、概率融合判决以及五状态机时序约束组成的级联框架，
从连续数据流中稳定提取车辆进入边界索引 $k_{\mathrm{in}}$ 与离开边界索引 $k_{\mathrm{out}}$，
输出车辆事件片段 $[k_{\mathrm{in}}, k_{\mathrm{out}})$ 作为后续处理的统一输入。
其中 $k_{\mathrm{out}}$ 表示恢复无车判决重新成立后的第一个采样点索引，因此事件片段采用半开区间 $[k_{\mathrm{in}}, k_{\mathrm{out}})$，
对应事件持续点数 $N_{\mathrm{occ}} = k_{\mathrm{out}} - k_{\mathrm{in}}$。

\paragraph*{(1) 三轴一阶差分特征与平滑处理}
设采样时刻 $k$ 的三轴地磁输出为
\begin{equation}
    \mathbf{B}(k) = \big[B_x(k),\,B_y(k),\,B_z(k)\big]^T .
\end{equation}
为抑制慢变基线并突出车辆扰动的快速变化分量，对三轴序列进行一阶差分：
\begin{equation}
    D_l(k) = B_l(k) - B_l(k-1),\quad l \in \{x, y, z\}.
\end{equation}
差分序列 $D_l(k)$ 的单位与 $B_l(k)$ 一致，仍为 \MagUnit。

差分会放大高频噪声与偶发尖峰，因此工程实现中需对差分序列进行平滑处理。
考虑端侧 MCU 的算力与实时性约束，本文采用 $L_s=\SmoothLen$ 点滑动平均：
\begin{equation}
    \overline{D}_l(k)=\frac{1}{L_s}\sum_{m=0}^{L_s-1} D_l(k-m),\quad l \in \{x, y, z\}.
\end{equation}
滑动平均可视为线性相位 FIR 滤波器的特例，其群时延为 $(L_s-1)/(2f_s)\approx 0.06\,\mathrm{s}$，
对边界定位影响可忽略。差分与平滑处理对三轴序列的作用关系示例如图~\ref{fig:diff_smooth} 所示。

\begin{figure}[!htbp]
    \centering
    \includegraphics[width=0.86\linewidth]{fig_diff_smooth}
    \caption{三轴差分特征与平滑处理示例（幅值单位为 \MagUnit）}
    \label{fig:diff_smooth}
\end{figure}

\paragraph*{(2) 背景统计建模与单轴概率得分}
将车辆检测建模为二元假设判决问题。定义时刻 $k$ 的车辆真实状态为 $e(k) \in \{0, 1\}$，
其中 $e(k) = 0$ 表示无车，$e(k) = 1$ 表示有车。
在无车状态下，平滑差分主要由测量噪声与偶发微扰构成，可用高斯近似描述：
\begin{equation}
    \overline{D}_l(k)\,|\,e(k) = 0 \sim \mathcal{N}\!\left(\mu_{l,0}, \sigma_{l,0}^2\right),\quad l \in \{x, y, z\}.
\end{equation}

工程实现中，选取连续无车时间窗 $\mathcal{K}_0$ 估计背景均值与尺度。为保证可复现，
本文固定背景窗口长度为 $|\mathcal{K}_0| = N_0 = \BgWinPts$（约 $10\,\mathrm{s}$ 的无车样本）：
\begin{align}
    \hat{\mu}_{l,0} &= \frac{1}{|\mathcal{K}_0|}\sum_{k \in \mathcal{K}_0}\overline{D}_l(k), \\
    \hat{\sigma}_{l,0} &= \sqrt{\frac{1}{|\mathcal{K}_0| - 1}\sum_{k \in \mathcal{K}_0}\left(\overline{D}_l(k) - \hat{\mu}_{l,0}\right)^2}.
\end{align}
为避免车辆扰动污染背景统计量，背景窗口仅在无车状态 $S_1$ 且 $P_{car}(k)<\theta_{arr}$ 时滚动更新；
一旦进入 $S_2\sim S_5$，冻结背景参数直至回到 $S_1$。

基于上述背景模型，定义单轴背景一致性得分 $p_{0,l}(k)$ 为观测值在无车分布下的双侧尾部概率：
\begin{equation}
    p_{0,l}(k) = 2\left[1 - \Phi\!\left(\frac{\left|\overline{D}_l(k)-\hat{\mu}_{l,0}\right|}{\hat{\sigma}_{l,0}}\right)\right],
\end{equation}
其中 $\Phi(\cdot)$ 为标准正态分布累积分布函数。为避免三轴乘积导致数值下溢，本文对得分作下界截断：
\begin{equation}
    p_{0,l}(k)\leftarrow \max\big(p_{0,l}(k),\,p_{\min}\big),\quad p_{\min}=\ProbFloor.
\end{equation}
为提高数值稳定性，工程实现亦可在对数域累计乘积项并在末端归一化；本文采用下界截断方案以保持实现简洁。
相应地，定义单轴扰动得分为
\begin{equation}
    p_{1,l}(k) = 1 - p_{0,l}(k).
\end{equation}

\paragraph*{(3) 三轴朴素融合与融合概率构造}
为避免直接拟合复杂的三轴联合分布，本文采用朴素融合思想对三轴得分进行合成，并引入相关性修正系数 $\rho_0, \rho_1$：
\begin{equation}
    \widetilde{P}_{env}(k) = \rho_0\prod_{l \in \{x, y, z\}}p_{0,l}(k),\qquad
    \widetilde{P}_{car}(k) = \rho_1\prod_{l \in \{x, y, z\}}p_{1,l}(k).
\end{equation}
在本文实验中，为保证实现简洁且可复现，取 $\rho_0=\rho_1=1$（朴素独立假设）。
进一步得到归一化融合概率：
\begin{align}
    P_{car}(k) &= \frac{\widetilde{P}_{car}(k)}{\widetilde{P}_{env}(k) + \widetilde{P}_{car}(k)}, \\
    P_{env}(k) &= 1 - P_{car}(k).
\end{align}
融合概率序列的典型形态及到达、离开阈值关系如图~\ref{fig:p_car_example} 所示。

\begin{figure}[!htbp]
    \centering
    \includegraphics[width=0.86\linewidth]{fig_p_car_example}
    \caption{融合概率序列示例及阈值示意（上半幅差分幅值单位为 \MagUnit）}
    \label{fig:p_car_example}
\end{figure}

\paragraph*{(4) 五状态有限状态机与事件片段输出（连续点数实现）}
为进一步抑制偶发脉冲干扰，并适应 \Fs\ 条件下波形时间尺度变化，
本文在融合概率序列 $P_{car}(k)$ 上引入五状态有限状态机进行时序判决。
状态机结构与状态转移条件如图~\ref{fig:fsm} 所示。定义状态空间为
\begin{equation}
    S = \{S_1, S_2, S_3, S_4, S_5\},
\end{equation}
其中 $S_1$ 为无车状态，$S_2$ 为到达检测状态，$S_3$ 为延时状态，$S_4$ 为车辆经过状态，$S_5$ 为离开检测状态。

考虑到嵌入式实现的计算开销与参数可控性，本文采用阈值与连续点数计数相结合的判决方式。
设到达与离开概率阈值分别为 $\theta_{arr}, \theta_{lea}$，连续计数阈值分别为 $N_{arr}, N_{lea}$，
防断裂延时点数为 $T_d$。定义到达与离开指示量：
\begin{align}
    I_{arr}(k) &= \mathbb{I}\!\left(P_{car}(k) \ge \theta_{arr}\right), \\
    I_{lea}(k) &= \mathbb{I}\!\left(P_{car}(k) \le \theta_{lea}\right),
\end{align}
并以递推方式维护连续计数器：
\begin{align}
    n_{arr}(k) &=
    \begin{cases}
        n_{arr}(k-1) + 1, & I_{arr}(k) = 1,\\
        0, & I_{arr}(k) = 0,
    \end{cases}
    \\
    n_{lea}(k) &=
    \begin{cases}
        n_{lea}(k-1) + 1, & I_{lea}(k) = 1,\\
        0, & I_{lea}(k) = 0,
    \end{cases}
\end{align}

\begin{figure}[!htbp]
    \centering
    \includegraphics[width=0.86\linewidth]{fig_fsm}
    \caption{五状态有限状态机示意}
    \label{fig:fsm}
\end{figure}

各状态的物理含义与跳变逻辑如下：
\begin{enumerate}
    \item \textbf{无车状态 $S_1$：}
        系统初始化后进入 $S_1$ 并持续监测 $P_{car}(k)$。
        当 $I_{arr}(k) = 1$ 时，转入 $S_2$ 并启动到达计数（令 $n_{arr} \leftarrow 1$）。

    \item \textbf{到达检测状态 $S_2$：}
        若后续连续满足 $I_{arr}(k) = 1$，则持续累加 $n_{arr}$；一旦出现 $I_{arr}(k) = 0$，回退至 $S_1$ 并清零 $n_{arr}$。
        当 $n_{arr}(k) \ge N_{arr}$ 时确认车辆到达，并定义车辆进入边界索引为
        \begin{equation}
            k_{\mathrm{in}} = k - N_{arr} + 1.
        \end{equation}
        确认到达后转入延时状态 $S_3$，并令延时计数 $n_d \leftarrow 0$。

    \item \textbf{延时状态 $S_3$：}
        车辆事件内部可能因车底弱磁区域或姿态变化出现短时回落。若到达确认后立即进入离开检测，
        该回落可能造成单车事件被误分割甚至重复计数。为此引入强制延时机制。
        在 $S_3$ 中递增延时计数器
        \begin{equation}
            n_d \leftarrow n_d + 1,\quad 0 \le n_d < T_d.
        \end{equation}
        当 $n_d \ge T_d$ 时转入 $S_4$。

    \item \textbf{车辆经过状态 $S_4$：}
        若 $I_{lea}(k) = 0$，保持在 $S_4$；
        当 $I_{lea}(k) = 1$ 时，转入 $S_5$ 并启动离开计数（令 $n_{lea} \leftarrow 1$）。

    \item \textbf{离开检测状态 $S_5$：}
        若后续连续满足 $I_{lea}(k) = 1$，则持续累加 $n_{lea}$；若出现 $I_{lea}(k) = 0$，回退至 $S_4$ 并清零 $n_{lea}$。
        当 $n_{lea}(k) \ge N_{lea}$ 时确认车辆离开，并定义车辆离开边界索引为
        \begin{equation}
            k_{\mathrm{out}} = k - N_{lea} + 1.
        \end{equation}
        其中 $k_{\mathrm{out}}$ 对应无车判决重新成立的起始索引，因此输出车辆事件片段 $[k_{\mathrm{in}}, k_{\mathrm{out}})$ 并返回 $S_1$。
\end{enumerate}

\paragraph*{(5) 参数设置原则与取值依据（\Fs，连续点数口径）}
为保证车辆检测算法在 $f_s = \Fs$ 条件下具备可复现性与工程鲁棒性，本文对关键参数给出统一的确定依据。
阈值与点数参数依据示例如图~\ref{fig:param_basis} 所示。
其中阈值满足 $\theta_{arr} > \theta_{lea}$；
点数阈值对应的确认时延可由 $N_{arr}/f_s$、$N_{lea}/f_s$ 与 $T_d/f_s$ 表示。

除表~\ref{tab:det_param} 中给出的阈值与点数参数外，本章进一步固化平滑滤波器、融合修正系数与背景窗口长度：
(1) 平滑采用 $L_s = \SmoothLen$ 点滑动平均；
(2) 取 $\rho_0 = \rho_1 = 1$；
(3) 背景窗口长度固定 $N_0 = \BgWinPts$，且仅在 $S_1$ 无车状态更新。

\begin{figure}[!htbp]
    \centering
    \includegraphics[width=0.86\linewidth]{fig_param_basis}
    \caption{关键参数取值的统计依据示例}
    \label{fig:param_basis}
\end{figure}

综合统计规律与工程约束，本文在 $f_s = \Fs$ 下采用的核心参数如表~\ref{tab:det_param} 所示。

\begin{table}[!htbp]
    \centering
    \caption{车辆检测关键参数设置（$f_s=\Fs$）}
    \label{tab:det_param}
    \begin{tabular}{c c c p{6.4cm}}
        \hline
        参数 & 符号 & 取值 & 取值依据（简述）\\
        \hline
        到达概率阈值 & $\theta_{arr}$ & 0.9 & 有车/无车 $P_{car}$ 分布分离，高阈值抑制虚警\\
        离开概率阈值 & $\theta_{lea}$ & 0.5 & 拖尾容错与离开确认稳定性折中，避免过早截断\\
        到达连续计数阈值 & $N_{arr}$ & 10 & 抑制无车段短触发；确认时延 $N_{arr}/f_s$\\
        离开连续计数阈值 & $N_{lea}$ & 10 & 抑制回弹误切分；确认时延 $N_{lea}/f_s$\\
        防断裂延时点数 & $T_d$ & 8 & 覆盖典型短时回落段；延时 $T_d/f_s$\\
        \hline
        平滑窗口长度 & $L_s$ & \SmoothLen & 端侧低算力下的噪声抑制；群时延约 $0.06$\,s\\
        融合修正系数 & $\rho_0,\rho_1$ & 1, 1 & 朴素独立假设，避免额外标定\\
        背景窗口长度 & $N_0$ & \BgWinPts & 约 $10$\,s 无车样本；仅在 $S_1$ 更新\\
        概率下界 & $p_{\min}$ & \ProbFloor & 防止乘积下溢与极端化\\
        \hline
    \end{tabular}
\end{table}

\FloatBarrier

\section{车辆信息检测平台}\label{sec:ch2_platform}

为实现路侧感知数据的可靠汇聚、远程管理与可视化应用，本文构建了由地磁检测节点、通信基站（网关）、云平台与可视化前端组成的车辆信息检测平台。
平台总体架构如图~\ref{fig:platform_arch} 所示。
平台在道路沿线按车道方向部署地磁检测节点，节点间距根据道路结构、空间覆盖与通信负载综合确定，一般取 $8 \sim 15\,\mathrm{m}$。
各节点具备独立设备标识与部署属性（如车道编号、安装位置等），在边缘侧完成第~\ref{sec:vehicle_detection}~节所述车辆检测与事件分割，
并以事件级摘要形式上报，实现边缘侧实时检测与云端统一管理的协同工作模式。

\begin{figure}[!htbp]
    \centering
    \includegraphics[width=0.88\linewidth]{fig_platform_arch}
    \caption{车辆信息检测平台总体架构示意}
    \label{fig:platform_arch}
\end{figure}

\paragraph{事件级数据流程与摘要字段}
在数据流程方面，地磁检测节点以采样率 $f_s = \Fs$ 采集三轴磁场序列，并在端侧实时输出车辆事件片段 $[k_{\mathrm{in}},\,k_{\mathrm{out}})$。
其中 $k_{\mathrm{in}}, k_{\mathrm{out}} \in \mathbb{Z}$ 表示事件在采样序列中的起止索引，且 $k_{\mathrm{out}}$ 为恢复无车判决重新成立后的第一个采样点索引。
事件持续点数定义为
\begin{equation}
    N_{\mathrm{occ}} = k_{\mathrm{out}} - k_{\mathrm{in}},
\end{equation}
相应驻留时长为
\begin{equation}
    T_{\mathrm{occ}} = \frac{N_{\mathrm{occ}}}{f_s}.
\end{equation}

为降低链路占用与节点能耗，系统常态运行仅上报事件级摘要，不回传大规模原始波形数据。
本文后续章节所需的三轴原始事件片段在实验采集阶段通过离线方式获取，用于标注、建库与算法验证，从而避免对常态通信链路与节点能耗造成额外负担。

为保证后续章节数据口径一致，事件级摘要字段建议按表~\ref{tab:event_summary} 固化，并预留扩展字段以支持上层应用迭代，实现底座检测与上层任务的解耦。

\begin{table}[!htbp]
    \centering
    \caption{事件级摘要字段定义（常态上报）}
    \label{tab:event_summary}
    \begin{tabular}{c c c p{7.2cm}}
        \hline
        字段类别 & 字段名称 & 符号/单位 & 说明 \\
        \hline
        标识类 & 设备标识 & $\mathrm{ID}_{dev}$ & 节点唯一标识 \\
        标识类 & 车道编号 & $\mathrm{ID}_{lane}$ & 节点所属车道或泊位编号 \\
        标识类 & 部署位置 & $\mathrm{ID}_{pos}$ & 安装位置或路段位置索引（可选） \\
        \hline
        时间类 & 事件时间戳 & $\tau$（s 或 ms） & 事件发生时间戳（可选，可定义为进入时刻时间戳） \\
        时间类 & 起止索引 & $k_{\mathrm{in}}, k_{\mathrm{out}}$ & 端侧检测得到的事件边界索引 \\
        统计类 & 驻留点数 & $N_{\mathrm{occ}}$ & $N_{\mathrm{occ}} = k_{\mathrm{out}} - k_{\mathrm{in}}$ \\
        统计类 & 驻留时长 & $T_{\mathrm{occ}}$（s） & $T_{\mathrm{occ}} = N_{\mathrm{occ}}/f_s$ \\
        \hline
        扩展类 & 质量标记 & flag & 异常、丢包、低电量等状态标记（可选） \\
        扩展类 & 业务标签 & label & 车型类别或停车占用状态等扩展字段（可选，后续章节使用） \\
        \hline
    \end{tabular}
\end{table}

\paragraph{通信基站与回传链路}
通信基站作为路侧汇聚网关，用于接收多个检测节点经 LoRa 链路上传的事件摘要数据，并通过蜂窝网络（如 4G）或有线网络回传至云平台。
考虑道路环境下无线链路的波动性，基站侧可配置缓存与重传机制，并结合去重与顺序恢复策略，提高端到端传输的连续性与可靠性。

\paragraph{云平台与前端可视化}
云平台承担事件数据的统一接入、解析校验、持久化存储与服务接口发布等任务。
平台将来自不同节点的事件记录按路段与车道维度进行组织，形成结构化数据集，用于支撑设备在线监测、事件检索与回放、
路段统计以及异常告警等功能。

平台前端提供面向运维与管理的图形化界面，用于展示道路部署概况、节点运行状态、事件记录与统计结果，其示例效果如图~\ref{fig:platform_ui} 所示。

\begin{figure}[!htbp]
    \centering
    \includegraphics[width=0.88\linewidth]{fig_platform_ui}
    \caption{车辆信息检测平台前端界面示例}
    \label{fig:platform_ui}
\end{figure}

综上，本文平台实现了边缘侧事件生成、基站汇聚转发、云端存储管理与前端可视化呈现的闭环链路，
为后续章节的数据组织、任务建模与实验验证提供支撑环境。

\FloatBarrier

\section{地磁车辆检测模块}\label{sec:ch2_module}

\paragraph{模块组成与选型说明}
为实现道路与停车场景下车辆信息的连续采集与低功耗上报，本文采用地磁车辆检测模块作为路侧嵌入式感知节点。
该模块面向长期在线部署，集成地磁感知、边缘计算、无线通信、供电与状态监测等功能单元，
能够在端侧完成车辆事件检测并将事件级摘要上传至上位平台。
本项目采用的硬件平台为既有工程选型：三轴地磁传感器选用 VCM5883，通信链路采用 LoRa，主控为低功耗 MCU，
并配套温度/电量监测与外部存储用于长期运行与断网缓存。
本文重点围绕事件级数据口径与算法流程开展研究，器件型号与功耗细节不在本章展开，但保留选型逻辑与功能分工说明。

\begin{table}[!htbp]
    \centering
    \caption{地磁检测节点主要硬件单元与功能分工（最低限度工程说明）}
    \label{tab:node_units}
    \begin{tabular}{c c p{8.0cm}}
        \hline
        功能单元 & 典型器件/接口 & 说明 \\
        \hline
        地磁感知 & VCM5883（三轴） & 输出三轴磁场分量 $\mathbf{B}(k)$，满足低功耗长期采样需求 \\
        边缘计算 & 低功耗 MCU & 完成采样、差分平滑、概率融合与状态机判决，输出事件边界与摘要字段 \\
        无线通信 & LoRa 模组 & 事件级摘要上报；支持多节点汇聚与低功耗长距离通信 \\
        状态监测 & 温度/电量检测 & 支持温漂相关监测与低电量告警，提高可维护性 \\
        本地存储 & Flash/EEPROM & 参数配置、断网缓存、运维数据记录等 \\
        \hline
    \end{tabular}
\end{table}

硬件实物如图~\ref{fig:module_photo} 所示。

\begin{figure}[!htbp]
    \centering
    \includegraphics[width=0.86\linewidth]{fig_module_photo}
    \caption{地磁车辆检测模块硬件实物图}
    \label{fig:module_photo}
\end{figure}

\paragraph{工作流程与数据口径}
地磁观测在三维空间中为矢量量。设采样时刻为 $k$，三轴地磁输出记为
\begin{equation}
    \mathbf{B}(k) = \big[B_x(k),\,B_y(k),\,B_z(k)\big]^T,
\end{equation}
其中 $B_x(k)$、$B_y(k)$、$B_z(k)$ 分别表示传感器在三个正交轴向的磁感应强度分量，单位为 \MagUnit。
本文采用的坐标系约定与第~\ref{sec:ch2_basic}~节一致；若传感器安装存在偏角，可通过标定完成轴系转换，从而保证端侧算法输入口径一致。

考虑到边缘节点的低功耗与长续航约束，本文采用 $f_s = \Fs$ 的采样模式，并与第~\ref{sec:vehicle_detection}~节算法参数保持一致。
模块工作流程可概括为：
\begin{enumerate}
    \item 节点按采样率采集 $\mathbf{B}(k)$，完成必要的预处理；
    \item 在端侧执行车辆检测与事件分割，输出车辆进入与离开边界索引 $k_{\mathrm{in}}$、$k_{\mathrm{out}}$，并形成车辆事件片段 $[k_{\mathrm{in}},\,k_{\mathrm{out}})$；
    \item 将事件级摘要按数据帧格式封装，经无线链路上传至通信基站或网关完成汇聚，并进一步回传至云平台存储与展示。
\end{enumerate}
其中 $k_{\mathrm{out}}$ 表示恢复无车判决重新成立后的第一个采样点索引。系统常态运行仅上传事件级摘要，以降低链路占用与节点能耗。

上述端侧检测与低负载上报构成的闭环机制，为后续章节的车型分类与稳定点停车占用判定提供一致的事件级数据入口与实验基础。

\FloatBarrier

\section{本章小结}\label{sec:ch2_summary}

本章围绕地磁车辆检测的基本原理与系统框架开展论述。
首先，从车辆引起地磁扰动的物理现象出发，给出了三轴磁场观测量的表示与坐标系约定，
并引入等效磁偶极子模型解释扰动的空间局部性、距离衰减与三轴差异性。
其次，面向 $f_s = \Fs$ 的应用场景，构建了车辆检测与事件分割的端侧处理流程：
通过三轴一阶差分抑制慢变基线并突出扰动变化分量（图~\ref{fig:diff_smooth}），基于背景统计模型构造单轴概率并进行融合（图~\ref{fig:p_car_example}），
进一步结合五状态机的时序约束（图~\ref{fig:fsm}）实现车辆进入与离开边界索引 $k_{\mathrm{in}}$、$k_{\mathrm{out}}$ 的稳定提取，
输出标准化车辆事件片段 $[k_{\mathrm{in}},\,k_{\mathrm{out}})$。
同时，为保证方法可复现，本章明确给出了平滑滤波器、融合修正系数与背景窗口等实现参数。
最后，介绍了由检测节点、通信基站与云端平台构成的车辆信息检测系统，实现事件级摘要数据的汇聚、存储与可视化展示。
系统常态运行仅上报事件级摘要字段，以降低链路占用与节点能耗；后续章节所需的三轴原始事件片段数据在实验采集阶段通过离线方式获取，用于标注与算法验证。
后续章节将以本章输出的车辆事件片段及其事件级摘要作为统一输入，开展方法设计与评估。








\chapter{基于单地磁传感器的车型分类方法}


\section{引言}

在智慧交通系统中，车流量与车型信息是交通运行状态评估、精细化管控以及交通规划的重要基础。
相较于仅进行车辆计数，车型信息能够进一步刻画交通组成结构，为差异化管控策略制定、道路资源配置与运行维护提供更具针对性的依据。
与此同时，实际道路环境具有车流密度高、工况变化快等特点，要求感知系统以较低成本实现长期稳定运行，并持续输出工程可用的车辆事件与属性信息。

地磁传感器通过感知车辆铁磁材料对地球磁场的扰动实现车辆检测，具有对光照和天气变化敏感性较低、节点可隐蔽部署、适于规模化布设等优势。
在工程部署中，采样频率、通信带宽与端侧算力共同决定系统成本与能耗水平。
本文所研究系统采用三轴地磁传感器进行采集，采样频率为 $50\,\mathrm{Hz}$，通信侧以事件级摘要上报为主。
在此条件下实现可靠的车型识别，有助于在不显著增加通信负担与能耗的前提下提升路侧节点输出信息的精细程度。

单节点单地磁传感器条件下的车型三分类具有显著工程难度。
一方面，单节点难以直接测得车速，车辆事件持续时间与车辆物理长度呈现耦合关系。
持续时间较短的事件既可能对应高速通过的大型车，也可能对应低速通过的小型车。
因此，依赖事件持续时间的简单判别方法在车速变化较大时容易失效，分类方法需要更多依赖波形形态信息与归一化处理。
另一方面，环境磁场随温度与地磁变化等因素产生缓慢漂移，若漂移抑制不足，车辆扰动的相对幅值与形态会发生偏移，影响特征一致性。
此外，交通数据中不同车型的样本数量分布往往不均衡，若不加处理，模型训练过程可能偏向样本数量较多的类别，导致少数类识别能力不足。

为应对上述问题，本章在第二章车辆事件检测底座的输出接口基础上，以单车事件片段作为分类输入，构建小型车、中型车与大型车三分类任务，并给出数据构建与标注口径。
在算法设计上，考虑端侧资源与通信约束，本章首先构建端侧轻量基线方法，围绕候选特征构造、特征筛选与低复杂度分类器实现可部署的车型识别，并通过类别不均衡处理策略缓解训练偏置。
进一步地，为减弱车速变化引起的时间伸缩影响，本章在离线条件下引入定长归一化与动态时间规整对齐，并构建一维卷积神经网络模型对对齐后的信号进行端到端学习，以评估更强模型在速度变化条件下的可达分类性能上限。
需要强调的是，所述深度模型用于离线仿真评估与对比分析，不作为端侧已部署方案。

本章后续内容安排如下：首先给出车型三分类任务定义、事件切片方式以及类别体系与标注规则；随后介绍端侧轻量基线的特征构造、特征筛选与分类器设计；接着阐述定长归一化、动态时间规整对齐与一维卷积网络的离线增强方案，并给出注意力机制与轻量化设计；最后通过对比实验与消融分析验证各模块对分类性能与鲁棒性的贡献，并讨论面向工程部署的资源与精度取舍。



\section{任务定义与数据构建}

本节给出车型三分类任务的定义、类别体系与标注口径，并说明数据采集、事件切片接口、预处理与数据集划分规则。
本节内容旨在形成统一数据口径，使端侧轻量基线与离线增强模型能够在一致条件下进行对比验证。

\subsection{类别体系与标注口径}
本研究采用三分类体系，类别分别为小型车、中型车与大型车。
类别划分以工程可执行性为原则，结合车辆外形尺寸、用途属性以及轴数等信息形成标注规则。
对存在歧义的车型设置合并原则，对无法确认类别或存在严重遮挡的样本进行剔除或单独记录，以保证训练与测试标签的一致性。
为提升论文可复现性，建议将三分类标注规则以表格形式给出，并在最终定稿时与实际标注口径保持一致。

\subsection{数据获取与事件切片接口}
地磁节点以固定采样频率输出三轴磁场分量序列。
第二章车辆事件检测底座在端侧完成车辆到达与离开时刻判定，并输出单车事件时间边界。
本章以该边界为切片依据，从连续磁场数据中截取单车事件片段，形成车型分类的基本样本单元。

设三轴原始磁场序列为
\[
\mathbf{B}[n]=\left[B_x[n],\,B_y[n],\,B_z[n]\right]^{\mathrm{T}},
\]
车辆事件片段对应的索引区间为 $n\in[n_{\mathrm{in}},\,n_{\mathrm{out}}]$。
事件片段长度记为 $N=n_{\mathrm{out}}-n_{\mathrm{in}}+1$。

系统在线链路以事件级摘要上报为主。
研发建库与离线评估阶段，可通过本地接口导出三轴波形片段用于训练与仿真评估。
该离线数据仅用于模型训练、参数分析与算法对比，不改变系统在线上报链路的基本约束。

\subsection{信号预处理与统一表示}
环境磁场存在缓慢漂移。
为削弱基线漂移对分类输入的影响，本章复用第二章底座模块提供的基线估计与事件切片结果，在事件片段上形成车辆扰动序列。
设事件到达前短窗口长度为 $N_0$，基线估计为
\[
\mathbf{B}_0=\frac{1}{N_0}\sum_{k=n_{\mathrm{in}}-N_0}^{n_{\mathrm{in}}-1}\mathbf{B}[k],
\]
车辆扰动序列表示为
\[
\Delta\mathbf{B}[n]=\mathbf{B}[n]-\mathbf{B}_0,\quad n\in[n_{\mathrm{in}},\,n_{\mathrm{out}}].
\]
后续端侧特征工程方法与离线深度模型均以 $\Delta\mathbf{B}[n]$ 作为统一输入表示。

由于传感器安装条件、车辆底盘高度及环境差异会造成扰动幅值变化，本章在离线训练与评估阶段对输入进行幅值归一化，以降低幅值尺度差异对分类器的影响。
幅值归一化方式可采用样本级归一化或基于训练集统计的标准化，具体形式与参数取值在实验设置中给出。

车辆速度变化会导致事件片段长度 $N$ 在样本间差异明显。
本章采用两类输入口径：对端侧轻量基线，允许输入为变长事件片段并提取统计特征；对离线深度模型，为满足批量训练输入维度一致性，将事件片段映射为定长序列，定长策略在第 3.4 节给出。

为便于复现与工程对齐，本章将关键预处理环节与输入表示口径汇总如表~\ref{tab:ch3_preprocess} 所示。

\begin{table}[htbp]
\centering
\caption{预处理与输入表示口径汇总}
\label{tab:ch3_preprocess}
\begin{tabular}{p{3.0cm}p{9.5cm}}
\hline
环节 & 口径说明 \\
\hline
事件切片 & 复用第二章事件检测底座输出的单车事件边界 $[t_{\mathrm{in}},t_{\mathrm{out}}]$ \\
基线扣除 & 复用第二章底座的基线估计结果，形成扰动序列 $\Delta\mathbf{B}[n]$ \\
幅值归一化 & 为降低幅值尺度差异影响，对输入进行归一化处理，具体形式见实验设置 \\
长度统一 & 端侧基线使用变长事件片段的统计特征；离线模型使用定长序列表示 \\
\hline
\end{tabular}
\end{table}

\subsection{数据集划分与评价指标}
为避免信息泄漏，数据集划分以车辆为最小单位进行，保证同一车辆的事件片段不会同时出现在训练集与测试集中。
划分过程中对三类车型进行分层抽样，尽量保持各集合中类别比例一致。
若数据覆盖不同日期与不同时段，划分时需兼顾时段分布，以评估模型对工况变化的鲁棒性。

评价指标采用总体准确率、宏平均指标与混淆矩阵。
宏平均指标用于降低类别不均衡对总体准确率的掩盖效应，混淆矩阵用于分析类别间主要混淆模式。
具体指标定义与实验参数在第3.5节统一给出。


\section{端侧轻量车型分类基线方法}

\subsection{基线框架与端侧输入输出}
面向在线部署需求，本节构建端侧轻量车型分类基线方法。
该方法以第二章事件检测底座输出的单车事件片段为输入，在端侧完成特征提取与轻量分类决策，并输出车型标签及置信度等摘要信息。
考虑通信链路以事件级摘要上报为主，端侧基线需在不依赖波形上传的条件下完成判别，且应具备较低的计算复杂度与存储开销。

设输入为基线扣除后的三轴扰动序列 $\Delta\mathbf{B}[n]$，输出为三分类标签 $\hat{y}\in\{1,2,3\}$，分别对应小型车、中型车与大型车。
分类置信度由分类器输出概率得到，用于后续异常样本分析与结果可信度表征。

单节点条件下无法直接测得车速，事件持续时间与车辆物理长度呈现耦合关系。
因此，本节端侧基线不以事件持续时间作为唯一判别依据，而是通过幅值、能量、形态结构与比值类特征综合刻画车辆扰动特性，以提升对速度变化的适应能力。
速度相关影响将在第3.5节通过分组实验进行验证与讨论。

\subsection{候选特征构造}
轻量基线的核心在于构造能够刻画车辆体量差异且实现成本可控的特征集合。
结合单地磁车辆分类常见做法，本研究候选特征以事件片段的时域统计量与形态量为主，并辅以少量复杂度可控的非线性特征。
候选特征主要包含以下几类。

第一类为幅值相关特征，用于表征扰动强度，例如各轴最大值、最小值、峰峰值以及峰值位置统计。
第二类为能量相关特征，用于表征事件整体扰动量，例如各轴能量、平均能量以及正负半轴能量占比等比值特征。
第三类为形态结构特征，用于刻画峰谷组合与局部变化，例如峰谷数量统计、过零点位置统计、极值间距等。
第四类为复杂度与不规则性特征，例如近似熵，用于补充对局部波动形态的描述。

为保证端侧可实现性，候选特征计算尽量采用一次遍历即可完成的统计量，避免引入频谱变换等高复杂度运算。
候选特征的完整清单建议以表格形式给出，包含特征名称、物理含义与计算方式，作为端侧实现依据。

\subsection{特征筛选与去冗余}
候选特征维度较高时，直接输入分类器可能引入冗余并放大噪声，从而增加过拟合风险与端侧开销。
因此，本节对候选特征进行重要性评估与去冗余筛选，形成低维稳定特征子集。

本研究采用三阶段筛选策略。
首先，基于树模型对特征重要性进行排序，得到候选特征的重要性序列。
其次，计算特征两两相关性，对相关性较高的特征对中重要性较低者进行剔除，以降低冗余。
最后，在剩余特征集合上进行逐步降维或遍历式选择，选取在验证集上性能与开销综合最优的特征子集。

相关性可采用皮尔逊相关系数定义：
\[
\rho_{ij}=\frac{\sum_{m}(x_{i,m}-\bar{x}_i)(x_{j,m}-\bar{x}_j)}
{\sqrt{\sum_{m}(x_{i,m}-\bar{x}_i)^2}\sqrt{\sum_{m}(x_{j,m}-\bar{x}_j)^2}}。
\]
相关性阈值与最终特征维度在实验设置中给出。
最终选择的特征子集建议以表格列出，包含特征名称、物理含义与计算方式，便于复现与端侧实现。

\subsection{轻量分类器与类别不均衡处理}
分类器选择以低复杂度模型为主，包括决策树、支持向量机与 Softmax 回归等。
其中 Softmax 回归能够输出类别概率分布，便于作为端侧置信度输出；支持向量机在中小规模特征空间上具有较好的判别能力；决策树具备一定可解释性，但需注意过拟合风险。

当三类样本数量分布不均衡时，标准交叉熵损失在训练过程中容易偏向样本数量较多的类别，从而导致少数类召回率下降。
为缓解该问题，本章在训练阶段引入类别加权机制。
设样本真实标签的独热向量为 $\mathbf{y}$，预测概率为 $\mathbf{p}$，类别权重为 $w_k$，则加权交叉熵可写为
\[
\mathcal{L}=-\sum_{k=1}^{3} w_k y_k \log p_k 。
\]
权重设置方式与具体取值在第3.5节给出，并结合宏平均指标与混淆矩阵评估其实际收益。

\subsection{端侧推理流程与复杂度分析}
端侧推理流程包括事件片段接收、基线扣除、特征提取、特征归一化与分类决策。
特征提取主要为线性遍历操作，其时间复杂度与事件长度近似呈线性关系。
分类决策阶段为低维向量运算，计算开销相对稳定。

端侧存储开销主要来自特征归一化参数与分类器参数。
在满足分类性能要求的前提下，应优先控制特征维度与模型参数规模，以降低节点能耗与内存占用。
本节讨论为端侧部署可行性提供复杂度口径，具体参数量与推理时间在第3.5节给出。




\section{时间对齐与一维卷积网络离线增强方法}

\subsection{方法总体框架与适用范围}
上一节给出了端侧轻量基线方法，其优势在于实现成本低、可部署性强。
然而，在车速变化较大或样本形态差异更复杂的条件下，基于固定尺度统计特征的方法可能面临判别边界不稳定的问题。
为评估更强模型在车速变化条件下的可达性能上限，并为后续模型压缩与迁移提供依据，本节在离线条件下构建基于时间对齐与一维卷积网络的增强方法。

需要强调的是，本文系统在线链路以事件级摘要上报为主。
本节方法依赖波形级输入与一定计算量，因此定位为离线仿真评估方法，不作为端侧在线已部署方案。
其输出结论用于对比分析与工程权衡讨论，并在第五章展望中为后续边云协同与模型迁移提供依据。

\subsection{定长归一化与动态时间规整对齐}
一维卷积网络训练通常要求输入维度一致，因此需将不同长度的事件片段映射为统一长度表示。
设三轴扰动序列为 $\Delta\mathbf{B}[n]$，长度为 $N$，目标长度为 $L$。
当 $N$ 大于 $L$ 时，采用插值方式对序列进行时间轴缩放，以尽量保留整体形态趋势。
当 $N$ 小于 $L$ 时，采用尾部补零方式扩展序列长度，避免直接截断造成的关键信息缺失。
目标长度 $L$ 的取值依据样本长度分布确定，以在信息保留与计算复杂度之间取得平衡，具体取值在第3.5节给出。

定长归一化本质上属于全局线性时间缩放。
当车辆通过过程存在局部加减速，或信号关键结构在不同样本间出现局部错位时，仅依靠线性缩放可能仍不足以实现形态对齐。
因此，本节进一步引入动态时间规整作为可选对齐模块。
动态时间规整允许时间轴进行非线性伸缩，通过寻找满足单调性与连续性约束的配准路径实现序列对齐，其原理与约束已在第3.3节给出。
本研究将动态时间规整限定为离线预处理，用于生成对齐后的定长输入，并在第3.5节通过消融实验评估其贡献。

\subsection{一维卷积网络结构设计与选型依据}
地磁扰动序列属于多通道时间序列信号。
在时间序列分类任务中，常见模型包括循环神经网络、基于自注意力的序列模型以及一维卷积网络。
本研究选择一维卷积网络作为离线增强模型的骨干结构，主要基于以下考虑。

第一，一维卷积网络依托局部感受野对局部形态结构具有较强建模能力，适合提取车辆磁扰动中的局部波包与峰谷组合特征。
第二，一维卷积网络具有较好的并行计算特性，训练与推理效率较高，有利于在后续轻量化与迁移分析中进行复杂度评估。
第三，在本文采样频率与事件长度条件下，样本有效长度处于中短序列范围，采用过于复杂的长序列建模结构可能带来更高的训练不确定性与过拟合风险。

本文构建由若干卷积块堆叠的一维卷积网络。
每个卷积块由一维卷积、归一化层、非线性激活与下采样组成。
随着网络深度增加，通道数逐步提升以增强表征能力，同时通过下采样扩大感受野。
在网络尾部采用全局平均池化得到定长特征向量，再通过全连接层输出三分类概率。
为保证可复现性，建议在正文中以表格列出网络各层的卷积核长度、通道数与下采样设置，并在实验设置中给出训练轮数、学习率与正则化策略。

\subsection{注意力机制与特征融合策略}
车辆通过事件中，不同轴向扰动贡献存在差异，且判别信息通常集中在车辆靠近传感器的关键时刻。
为增强模型对关键通道与关键时间片段的关注能力，可在卷积网络中引入注意力机制，对中间特征进行加权重标定。

通道注意力用于学习不同通道的重要性权重，其物理含义对应三轴分量在不同车辆与不同工况下的贡献差异。
时间注意力用于学习时间位置权重，其物理含义对应车辆关键部件经过传感器附近时产生的高信息量波段。
注意力模块会带来额外计算开销，因此需要通过消融实验评估其增益与代价，并在轻量化设计中进行取舍。

在样本规模有限时，纯端到端模型可能出现过拟合或收敛不稳定。
为增强模型可解释性与样本效率，本研究进一步尝试将卷积网络学习到的深层特征与少量物理含义明确的手工特征进行融合。
融合方式可在网络尾部对深层特征与手工特征进行拼接，再通过分类头输出三分类概率。
该策略的有效性与适用条件将在第3.5节通过对比实验给出。

\subsection{模型轻量化与训练设置}
考虑到后续潜在的端侧迁移需求，本研究在离线仿真中设计并评估一组轻量化结构变体。
轻量化主要包括减少网络层数、降低通道数，以及减少注意力模块插入位置或简化注意力结构。
轻量化可显著降低参数规模与推理开销，但可能带来精度下降，需结合混淆矩阵与宏平均指标进行综合权衡。

离线训练采用交叉熵损失，并可在类别不均衡条件下引入类别加权机制以缓解训练偏置。
为提升泛化能力，可采用权重衰减、随机失活与早停等策略，并在不改变车辆磁扰动基本形态的前提下进行适度数据增强，例如幅值缩放与加性噪声扰动。
具体超参数与训练策略在第3.5节给出。


\section{实验测试结果与性能分析}

本节在统一的数据构建与评价口径下，对端侧轻量基线方法与离线增强模型进行对比验证。
其中，端侧方法面向在线部署需求，离线增强方法用于评估在车速变化条件下的可达分类性能上限，并为后续模型轻量化与工程迁移提供参考。

\subsection{实验数据与设置}
\subsubsection{数据集统计与划分}
实验数据来源于真实道路采集。
以第二章事件检测底座输出的单车事件区间为切片依据，形成车型三分类样本。
数据集类别构成与训练集、验证集、测试集划分见表~\ref{tab:ch3_dataset}。
为避免信息泄漏，划分以车辆为最小单位进行，保证同一车辆的事件片段不会同时出现在训练集与测试集中。

\begin{table}[htbp]
\centering
\caption{车型三分类数据集统计与划分}
\label{tab:ch3_dataset}
\begin{tabular}{lcccc}
\hline
类别 & 总样本数 & 训练集 & 验证集 & 测试集 \\
\hline
小型车 & 待补充 & 待补充 & 待补充 & 待补充 \\
中型车 & 待补充 & 待补充 & 待补充 & 待补充 \\
大型车 & 待补充 & 待补充 & 待补充 & 待补充 \\
\hline
合计 & 待补充 & 待补充 & 待补充 & 待补充 \\
\hline
\end{tabular}
\end{table}

\subsubsection{对比方法与实现细节}
对比方法包括端侧轻量基线、离线卷积网络基线、离线增强模型及其消融变体。
端侧基线的特征维度与分类器参数、离线模型的输入长度、网络结构与训练设置在此处给出，并在全组实验中保持一致。
若采用类别加权机制，则需在此明确权重设置方式与具体取值。

\subsubsection{评价指标}
本章采用总体准确率作为主要指标，并给出宏平均指标与混淆矩阵以反映类别不均衡条件下的真实性能。
宏平均指标定义如下：
\[
F_{1,k}=\frac{2P_kR_k}{P_k+R_k},
\qquad
F_{1,\mathrm{macro}}=\frac{1}{3}\sum_{k=1}^{3}F_{1,k}.
\]

\subsection{总体性能对比}
表~\ref{tab:ch3_overall} 汇总端侧基线与离线模型的总体性能指标。
混淆矩阵如图~\ref{fig:ch3_cm_baseline} 与图~\ref{fig:ch3_cm_dl} 所示。
在最终结果填入后，本节将结合混淆矩阵解释中型车与相邻类别的主要混淆模式，并分析其可能原因。

\begin{table}[htbp]
\centering
\caption{总体性能对比}
\label{tab:ch3_overall}
\begin{tabular}{lccc}
\hline
方法 & 总体准确率 & 宏平均 $F_1$ & 备注 \\
\hline
端侧轻量基线 & 待补充 & 待补充 & 特征工程与轻量分类器 \\
离线卷积网络基线 & 待补充 & 待补充 & 定长输入与端到端学习 \\
离线增强模型 & 待补充 & 待补充 & 时间对齐与注意力等 \\
\hline
\end{tabular}
\end{table}

% \begin{figure}[htbp]
% \centering
% \includegraphics[width=0.70\linewidth]{fig_ch3_cm_baseline}
% \caption{端侧轻量基线混淆矩阵}
% \label{fig:ch3_cm_baseline}
% \end{figure}

% \begin{figure}[htbp]
% \centering
% \includegraphics[width=0.70\linewidth]{fig_ch3_cm_dl}
% \caption{离线增强模型混淆矩阵}
% \label{fig:ch3_cm_dl}
% \end{figure}

\subsection{速度鲁棒性分析}
为验证车速伸缩对分类性能的影响，本节将样本按速度区间或事件持续时间分组，并分别统计各方法在不同分组上的性能。
若缺少直接速度标注，可采用事件持续时间分组作为速度代理，并在正文中说明该代理的合理性与局限性。
分组规则与统计结果见表~\ref{tab:ch3_speed_group}。

\begin{table}[htbp]
\centering
\caption{分组条件下的性能对比}
\label{tab:ch3_speed_group}
\begin{tabular}{lcccc}
\hline
分组 & 样本数 & 基线准确率 & 增强模型准确率 & 宏平均 $F_1$ \\
\hline
低速或长事件 & 待补充 & 待补充 & 待补充 & 待补充 \\
中速或中事件 & 待补充 & 待补充 & 待补充 & 待补充 \\
高速或短事件 & 待补充 & 待补充 & 待补充 & 待补充 \\
\hline
\end{tabular}
\end{table}

\subsection{消融实验}
为明确各模块贡献，本节设置两组关键消融。
第一组用于评估时间归一化与对齐策略对速度鲁棒性的贡献，结果见表~\ref{tab:ch3_ablation_align}。
第二组用于评估注意力模块的有效性与代价，结果见表~\ref{tab:ch3_ablation_att}。
在最终结果填入后，本节将结合宏平均指标与混淆矩阵讨论其对少数类的影响。

\begin{table}[htbp]
\centering
\caption{时间归一化与对齐策略消融}
\label{tab:ch3_ablation_align}
\begin{tabular}{lccc}
\hline
变体 & 输入策略 & 总体准确率 & 宏平均 $F_1$ \\
\hline
变体一 & 不进行对齐与定长 & 待补充 & 待补充 \\
变体二 & 定长归一化 & 待补充 & 待补充 \\
变体三 & 定长归一化并引入动态时间规整 & 待补充 & 待补充 \\
\hline
\end{tabular}
\end{table}

\begin{table}[htbp]
\centering
\caption{注意力模块消融}
\label{tab:ch3_ablation_att}
\begin{tabular}{lccc}
\hline
变体 & 注意力形式 & 总体准确率 & 宏平均 $F_1$ \\
\hline
变体一 & 不使用注意力 & 待补充 & 待补充 \\
变体二 & 通道注意力 & 待补充 & 待补充 \\
变体三 & 时间注意力 & 待补充 & 待补充 \\
变体四 & 通道与时间联合注意力 & 待补充 & 待补充 \\
\hline
\end{tabular}
\end{table}

\subsection{复杂度评估与部署讨论}
为评估模型向端侧迁移的潜在可行性，本节统计各模型变体的参数规模、模型存储占用与推理时间，并形成表~\ref{tab:ch3_complexity}。
推理时间可在通用处理器上测得，并给出端侧估算口径。
需要强调的是，本节讨论用于复杂度与精度的权衡分析，不构成端侧已部署深度模型的结论。

\begin{table}[htbp]
\centering
\caption{模型复杂度对比}
\label{tab:ch3_complexity}
\begin{tabular}{lcccc}
\hline
模型 & 参数量 & 模型大小 & 推理时间 & 备注 \\
\hline
端侧轻量基线 & 待补充 & 待补充 & 待补充 & 特征与轻量分类 \\
离线增强模型 & 待补充 & 待补充 & 待补充 & 卷积网络与注意力 \\
轻量化变体一 & 待补充 & 待补充 & 待补充 & 减层或减通道 \\
轻量化变体二 & 待补充 & 待补充 & 待补充 & 注意力删减等 \\
\hline
\end{tabular}
\end{table}

\subsection{误差来源与典型样本分析}
除总体指标外，有必要结合典型误判样本对误差来源进行解释。
本节重点分析中型车误差来源，原因在于中型车在体量与形态上处于两端类别之间的过渡区域，其磁扰动形态更易与小型车或大型车产生混叠。
结合混淆矩阵，可从以下角度开展归因分析。

第一，速度变化导致的时间伸缩会改变局部结构的相对位置，若输入归一化与对齐不足，模型可能将高速通过的中型车误判为小型车，或将低速通过的中型车误判为大型车。
第二，事件切片边界存在误差时，车辆头部或尾部信息可能被截断，影响持续时间与形态结构特征，进而加剧中型车与相邻类别的混淆。
第三，当交通流较密集时，相邻车辆扰动可能发生叠加，使得事件片段包含额外波动结构，从而影响模型对关键峰谷结构的识别。

建议在最终论文中选取若干中型车误判样本给出波形对比图，并与正确分类样本进行对照，形成现象与原因的对应关系，以增强结论解释力。



\section{本章小结}

本章在第二章车辆事件检测底座输出接口的基础上，构建了单节点三轴地磁条件下的车型三分类任务，并围绕工程部署约束给出了两条互补技术路线。

面向在线部署需求，本章构建端侧轻量基线方法，通过事件片段特征提取与轻量分类器实现车型识别，并针对类别不均衡问题引入相应的缓解策略，为事件级摘要上报链路下的车型识别提供了可部署方案。

面向离线仿真评估需求，本章进一步引入定长归一化与动态时间规整对齐，并构建一维卷积网络模型对波形级输入进行端到端学习，讨论注意力机制、特征融合与模型轻量化策略在提升分类性能与降低复杂度方面的作用。
通过对比实验、消融分析与复杂度评估，本章为后续模型压缩、端侧迁移与边云协同方案设计提供了方法依据与工程权衡口径。

下一章将面向道路停车与占用判定任务，围绕信号稳定点构建可部署的停车检测方法，并在复杂工况下验证其鲁棒性。































































\chapter{基于稳定点的道路停车与占用检测方法}
\label{chap:ch4_parking}

\section{引言}
\label{sec:ch4_intro}

公路场景下的异常停车与占用事件会显著降低道路通行效率，并可能诱发追尾及二次事故等安全风险，因此需要在事件发生后尽可能早地检测并上报，为交通管理与应急处置提供依据。与停车场车位占用监测相比，道路场景车流密度波动更大、邻道干扰更频繁，异常停车往往需要在连续车流背景下被可靠地区分。

本章研究的监测场景如图\ref{fig:ch4_scene}所示：以高速公路或城市快速路为代表的多车道道路中，地磁传感器节点埋设于应急车道或路肩停车带区域，用于监测该区域内车辆的停靠与占用状态。由于传感器通常布设在靠近车道分界线的位置，主车道车辆的持续通行会对节点产生短时或频繁的磁场扰动，使应急车道停车事件需要在连续车流背景下被识别，这对端侧单传感器算法的稳定性与鲁棒性提出了更高要求。

第2章构建了基于单地磁三轴信号（采样率为$50\,\mathrm{Hz}$）的车辆事件检测底座，通过一阶差分与状态机时序约束实现车辆扰动事件的稳定检测，并输出事件片段及摘要字段，为后续高层语义识别提供统一输入接口。本章以第2章事件级输出为输入，在端侧单传感器、低算力与低通信开销等约束下，进一步研究道路停车与占用检测问题，目标是在事件级输出的基础上识别停车确认、占用维持与释放边界等高层语义状态。

需要指出的是，停车与占用语义并不能由“事件触发--事件结束”直接刻画。当车辆停稳后，磁场变化趋缓，一阶差分回落至零附近，其形态与无车状态下的差分特征高度相似，导致仅依赖差分事件难以区分车辆通过与车辆在检测范围内停靠并形成占用。此外，在连续车流条件下，传感器附近可能长期处于扰动状态，使稳定窗难以出现；缓行或停走又会造成变化速率降低的伪稳态现象。若采用宽松的参考更新条件，容易发生参考误更新与误判。因此，停车确认与占用维持逻辑需要同时约束窗内波动与窗间水平差异，以避免将缓慢变化误认为稳定。

针对上述问题，本章提出一种基于地磁稳定点的停车与占用检测方法。该方法利用车辆停稳或无车通过时的稳定点动态表征局部磁场稳态，并以此构造动态参考以替代固定基线，从而提升对环境漂移与工况变化的适应性。在单车停靠场景下，通过比较停靠前后稳定点的漂移量实现停车确认；在占用持续场景下，通过漂移相似性度量区分“停靠车仍在”与“新车通过”；在连续车流导致稳定点不可得时，引入固定窗均值作为稳定点近似，并结合一致性计数器实现稳健确认。

本章主要工作包括三方面：其一，提出稳定点驱动的动态参考表示，实现停车与占用稳态水平的显式表征；其二，构造稳定点漂移与相似性判据，并结合外层状态机实现占用维持与释放的可解释判定；其三，针对连续车流稳定窗缺失，引入固定窗近似与一致性计数器，提高拥堵工况下的可用性与鲁棒性。

% \begin{figure}[t]
% \centering
% \includegraphics[width=0.92\linewidth]{figures/ch4_scene.pdf}
% \caption{道路停车与占用监测场景示意：传感器节点布设于应急车道或路肩停车带；主车道车辆连续通行产生邻道扰动；监测目标为停靠车辆的确认、占用维持与释放边界。}
% \label{fig:ch4_scene}
% \end{figure}

本章结构如下：第\ref{sec:ch4_problem}节给出停车与占用检测的任务定义、状态口径与复杂工况；第\ref{sec:ch4_method}节介绍稳定点提取、动态参考构建、漂移与相似性判据及外层状态机约束；第\ref{sec:ch4_robust}节给出漂移抑制、伪稳态抑制与连续车流处理策略；第\ref{sec:ch4_exp}节通过多工况数据验证方法有效性并进行对比与消融；第\ref{sec:ch4_summary}节总结本章工作。

\section{停车与占用问题定义与信号特性}
\label{sec:ch4_problem}

\subsection{任务定义与状态口径}
\label{subsec:ch4_def}

本章以第2章输出的车辆扰动事件片段为输入，面向道路侧停车与占用检测任务，输出面向管理侧的语义状态序列。为与端侧事件检测底座保持一致，本文将停车与占用过程抽象为以下状态与边界：
\begin{itemize}
  \item 空闲（\texttt{FREE}）：检测区域内无车辆占用，磁场处于环境稳态附近；
  \item 停车确认（\texttt{PARK\_CONFIRM}）：车辆到达后在检测范围内停稳并形成稳态偏置；
  \item 占用维持（\texttt{OCCUPIED}）：车辆持续停靠导致稳态偏置保持，期间可能叠加邻道扰动；
  \item 占用释放（\texttt{RELEASE}）：停靠车辆离开并回到环境稳态附近。
\end{itemize}
其中，停车确认与占用释放的边界分别对应占用状态的建立与解除。需要强调的是，\texttt{diff\_on}/\texttt{diff\_off}仅刻画扰动开始与结束，并不等价于占用建立与解除，因此需引入稳态水平表征与时序约束实现语义判定。

\subsection{停车过程的简化运动模型}
\label{subsec:ch4_motion}

为解释车辆由行驶到停稳过程中磁场信号由扰动向近似稳态演化的规律，本节给出停车过程的简化运动描述。建立道路坐标系：$x$为车行方向，$y$为横向，$z$为竖直方向。为简化分析，假设车辆沿$x$方向做直线运动，其余方向位移近似不变。将停车过程抽象为减速驶入与停止驻留两阶段：车辆以初速度$v_0>0$驶入检测区域，并以减速度幅值$a>0$匀减速直线运动直至停稳，则车辆纵向位移可表示为
\begin{equation}
x(t)=
\begin{cases}
x_0+v_0 t-\dfrac{1}{2}a t^2, & 0 \le t < t_s,\\[4pt]
x_0+\dfrac{v_0^2}{2a}, & t \ge t_s,
\end{cases}
\qquad
t_s=\dfrac{v_0}{a}.
\label{eq:ch4_motion}
\end{equation}
式\eqref{eq:ch4_motion}用于定性说明停车过程的时间尺度与位移变化关系，后续算法不依赖精确物理参数估计。

进一步地，速度会导致磁场波形在时间轴上的伸缩。记三轴磁场为$\mathbf{B}(t)=\mathbf{B}(x(t),y_0,z_0)$，则对时间求导可得
\begin{equation}
\dfrac{\mathrm{d}\mathbf{B}(t)}{\mathrm{d}t}
=
\dfrac{\partial \mathbf{B}}{\partial x}\, v(t),
\label{eq:ch4_speed}
\end{equation}
其中$v(t)=\mathrm{d}x(t)/\mathrm{d}t$为车辆瞬时速度。式\eqref{eq:ch4_speed}表明磁场变化速率与车辆速度成正比，因此在低速缓行或停走工况下，一阶差分幅值可能显著变小而呈现伪稳态现象。为避免将“变化慢”误判为“稳定”，后续稳定点与动态参考机制需要同时约束窗内波动与窗间水平差异。

\subsection{等效磁偶极子扰动模型与典型波形形态}
\label{subsec:ch4_dipole}

为在仿真层面生成三轴磁场的典型波形，本节给出等效磁偶极子扰动模型。将车辆在传感器附近产生的磁扰动近似为若干等效磁偶极子之和。设传感器位置为$\mathbf{p}_s$，第$i$个偶极子位置为$\mathbf{p}_i(t)$、磁矩为$\mathbf{m}_i$，则其在传感器处产生的磁场扰动为
\begin{equation}
\Delta\mathbf{B}_i(t)=
\dfrac{\mu_0}{4\pi}
\left(
\dfrac{3\mathbf{r}_i(t)\big(\mathbf{m}_i^{\mathsf{T}}\mathbf{r}_i(t)\big)}{\|\mathbf{r}_i(t)\|^5}
-
\dfrac{\mathbf{m}_i}{\|\mathbf{r}_i(t)\|^3}
\right),
\qquad
\mathbf{r}_i(t)=\mathbf{p}_s-\mathbf{p}_i(t),
\label{eq:ch4_dipole}
\end{equation}
其中$\mu_0$为真空磁导率，$\mathbf{r}_i(t)$为从偶极子指向传感器的相对位置向量。三轴传感器测得的磁场序列可表示为
\begin{equation}
\mathbf{B}(t)=\mathbf{B}_{\mathrm{env}} + \sum_{i=1}^{N}\Delta\mathbf{B}_i(t) + \boldsymbol{\eta}(t),
\label{eq:ch4_meas}
\end{equation}
其中$\mathbf{B}_{\mathrm{env}}$为环境稳态磁场，$\boldsymbol{\eta}(t)$包含测量噪声与低频漂移项。

由式\eqref{eq:ch4_motion}--\eqref{eq:ch4_meas}可定性得到两类典型波形：
\begin{itemize}
  \item 正常通过：车辆以近似恒速穿越敏感区，$\Delta\mathbf{B}(t)$呈现一次性扰动并在车辆远离后回归环境稳态附近；
  \item 减速停靠：车辆驶入后逐渐停稳，$\Delta\mathbf{B}(t)$在停稳后收敛至与车辆相对位置相关的稳态偏置，并在占用期间保持，期间可能叠加邻道车辆造成的短时扰动。
\end{itemize}
由于停稳阶段差分信号接近零，其特征与空闲状态高度相似，因此停车确认与占用维持需引入稳定点与动态参考机制进行语义区分。
\subsection{任务定义、状态口径与复杂工况}
\label{subsec:ch4_definition}

\paragraph{监测场景与问题边界}
本章面向道路旁侧停靠车道（应急车道或路肩）的异常停车与占用监测。传感器节点部署在停靠车道附近，
需在相邻主车道连续车流背景下，识别停靠车辆的建立、占用维持与释放边界。与停车位占用监测相比，
该场景的关键难点在于：相邻车道车辆持续经过会对三轴地磁序列产生弱扰动，导致稳定窗难以出现或出现延迟，
从而使停车确认与释放判定更易受干扰。

\paragraph{任务定义与输入输出接口}
第\ref{chap:ch2_event}章输出的事件检测底座以一阶差分触发与状态机时序约束为核心，给出车辆扰动事件片段。
记第 $m$ 次扰动事件为
\begin{equation}
\mathcal{E}_m = \bigl(t_{\mathrm{in},m},\, t_{\mathrm{out},m}\bigr),
\end{equation}
其中 $t_{\mathrm{in},m}$ 与 $t_{\mathrm{out},m}$ 分别对应 \texttt{diff\_on} 与 \texttt{diff\_off} 时刻，
并可获得事件邻域的预处理后三轴地磁序列 $\widetilde{\mathbf{B}}(k)$。

本章目标是在事件级输出之上形成可上报的停车与占用语义状态，输出包括：
停车确认标志 \texttt{park\_flag}、占用状态 \texttt{occ\_flag}，以及停车开始时刻 $t_{\mathrm{park,in}}$、
停车结束时刻 $t_{\mathrm{park,out}}$ 与停车持续时间 $T_{\mathrm{park}}$ 等字段。

\paragraph{状态口径与时间戳对齐原则}
停车与占用语义不能由 ``扰动开始--扰动结束'' 直接刻画。其原因在于：车辆停稳后磁场变化趋缓，
一阶差分回落至零附近，差分事件检测会将停靠过程仅解释为一次到达与一次离开，缺乏对停车语义的感知。:contentReference[oaicite:3]{index=3}
因此，本章采用稳定点与动态参考来刻画稳态水平，并在此基础上定义停车确认与占用维持逻辑。

为保持与第\ref{chap:ch2_event}章接口一致，本文将停车边界与事件片段末端对齐：
\begin{itemize}
  \item \textbf{停车开始边界}：若事件 $\mathcal{E}_{m}$ 被判定为 ``驶入并停靠''，则定义
  \begin{equation}
  t_{\mathrm{park,in}} \triangleq t_{\mathrm{out},m}.
  \end{equation}
  该定义表示车辆动态扰动结束并进入候选稳态段的起点；停车确认可能需要额外的稳定窗观测时间，
  因此算法输出的确认时刻可记为 $\hat{t}_{\mathrm{park,in}}$，并在实验中量化确认延迟。
  \item \textbf{停车结束边界}：占用维持过程中，车辆驶离会触发新的扰动事件 $\mathcal{E}_{n}$。
  若该事件被判定为 ``停靠车辆释放''，则定义
  \begin{equation}
  t_{\mathrm{park,out}} \triangleq t_{\mathrm{out},n}.
  \end{equation}
\end{itemize}
停车持续时间定义为
\begin{equation}
T_{\mathrm{park}} = t_{\mathrm{park,out}} - t_{\mathrm{park,in}}.
\label{eq:ch4_Tpark}
\end{equation}
在在线监测中，若尚未观测到释放边界，可用 $T_{\mathrm{park}}(t)=t-t_{\mathrm{park,in}}$ 进行滚动估计。

给定业务阈值 $T_{\mathrm{short}}$ 与 $T_{\mathrm{occ}}$（且 $T_{\mathrm{occ}} > T_{\mathrm{short}}$），
停车事件可按持续时长划分为短停与长停。需要强调的是：\texttt{occ\_flag} 用于表达当前是否处于占用持续态
（\texttt{OCCUPIED}），其置位发生在停车确认之后；短停与长停属于由持续时间阈值派生的业务判定，
可在在线上报或离线统计中由 $t_{\mathrm{park,in}}$ 与当前时刻（或 $t_{\mathrm{park,out}}$）计算得到。

\paragraph{稳定判据与复杂工况}
道路场景影响停车与占用判定的关键因素包括：
\begin{enumerate}
  \item \textbf{环境慢漂移与参考维护困难}：温度、供电与周边电磁环境变化会引起背景磁场缓慢漂移。
  传统固定基线或宽松更新策略容易发生参考误更新。稳定点方法通过动态反映无车或停稳时的真实稳态，
  用以替代固定基线并提升环境适应性。:contentReference[oaicite:4]{index=4}
  \item \textbf{连续车流导致稳定窗缺失}：相邻车道车辆频繁经过会使传感器长期处于扰动状态，
  稳定点难以直接获得，停车确认延迟增大甚至漏检。为此需要设计连续车流下的退化分支，
  例如采用固定时间窗均值近似稳定点并引入一致性计数器进行稳健确认。:contentReference[oaicite:5]{index=5}
  \item \textbf{伪稳态（缓行或停走）}：缓慢变化并不等价于稳定。仅用窗内波动幅度作为稳定判据，
  容易将缓行车辆导致的缓慢变化误判为稳态。为抑制该问题，需要同时约束窗内波动与窗间水平差异，
  其中窗内极差刻画波动范围，窗间均值差刻画相邻窗口磁场水平差，用以区分缓行造成的缓慢变化与环境缓慢漂移。:contentReference[oaicite:6]{index=6}
  \item \textbf{邻道弱扰动污染稳态统计}：邻道车辆产生的弱扰动可能不触发明显的差分事件，
  但会抬高稳定窗的波动统计量或引入水平偏移，因此需要在后续方法中设置严格的门控条件与状态转移约束。
\end{enumerate}

表\ref{tab:ch4_keyparams} 汇总本章核心符号与参数，后续方法与实验均沿用该口径。

\begin{table}[t]
\centering
\caption{本章核心符号与参数说明}
\label{tab:ch4_keyparams}
\begin{tabular}{c p{0.28\linewidth} p{0.60\linewidth}}
\hline
符号 & 含义 & 说明 \\
\hline
$f_s$ & 采样率 & 本文为 $50\,\mathrm{Hz}$ \\
$L$ & 稳定窗长度 & 滑动窗口长度（采样点数），对应时间为 $L/f_s$ \\
$s$ & 窗步长 & 默认取 $s=L$，用于比较相邻窗口统计量 \\
$R_{\mathrm{th}}$ & 波动阈值 & 约束窗内波动量（如极差或其融合量） \\
$M_{\mathrm{th}}$ & 水平差阈值 & 约束窗间水平差（如均值差或其融合量） \\
$D_{\mathrm{th}}$ & 停车漂移阈值 & 停车确认所需稳态漂移幅值下限 \\
$D_{\mathrm{free}}$ & 回归阈值 & 释放判定需回归环境参考的程度 \\
$\mathrm{dist}_{\mathrm{th}}$ & 相似性阈值 & 区分占用维持与释放的关键阈值（如欧氏距离门限） \\
$T_{\mathrm{seek}}$ & 寻稳超时 & 事件结束后搜索稳定点的最长等待时间 \\
$L_{\mathrm{fix}}$, $c_{\mathrm{th}}$ & 退化分支参数 & 固定窗长度与一致性计数器门限 \\
\hline
\end{tabular}
\end{table}

\section{基于稳定点的动态参考停车与占用判定方法}
\label{sec:ch4_method}

本节给出停车与占用判定的核心算法。总体思路为：以第2章事件片段为触发入口，在事件前后搜索稳定窗并提取稳定点，构建环境参考稳定点$\mathbf{S}_{\mathrm{pre}}$与占用参考稳定点$\mathbf{S}_{\mathrm{post}}$；通过稳定点漂移实现停车确认，通过相似性度量区分过车干扰与释放或占用状态变化；当连续车流导致稳定窗不可得时，启用固定窗近似与一致性计数器退化分支以保持可用性。

\subsection{总体流程与变量约定}
\label{subsec:ch4_vars}

设预处理后三轴地磁序列为
\begin{equation}
\widetilde{\mathbf{B}}(k) = \left[\widetilde{B}_x(k), \, \widetilde{B}_y(k), \, \widetilde{B}_z(k)\right]^{\mathsf{T}}, \qquad k \in \mathbb{Z},
\label{eq:ch4_Bvec}
\end{equation}
采样率为$f_s = 50\,\mathrm{Hz}$，采样点与时间满足$t=k/f_s$。

第2章事件检测底座输出第$m$次扰动事件片段
\begin{equation}
\mathcal{E}_m = \bigl(t_{\mathrm{in},m},\, t_{\mathrm{out},m}\bigr),
\end{equation}
其中$t_{\mathrm{in},m}$与$t_{\mathrm{out},m}$分别对应\texttt{diff\_on}与\texttt{diff\_off}时刻。

外层状态机输出包括：停车确认事件\texttt{park\_flag}（仅在确认瞬时置1并上报）、占用状态\texttt{occ\_flag}（占用期间维持为1）、以及$t_{\mathrm{park,in}}$、$t_{\mathrm{park,out}}$与$T_{\mathrm{park}}$。
内部变量包括：稳定状态$g_{\mathrm{st}}(k)$、稳定点触发标志$f_{\mathrm{st}}(k)$、环境参考$\mathbf{S}_{\mathrm{pre}}$、占用参考$\mathbf{S}_{\mathrm{post}}$、占用态事件后的新稳定点$\mathbf{S}_{\mathrm{new}}$、漂移向量$\Delta\mathbf{B}$及其幅值$\Delta B$、相似性度量$\mathrm{dist}$、以及连续车流退化分支中的一致性计数器$c_n$等。

算法以事件片段为触发入口：在事件前后限定的时间区间内搜索稳定窗并锁定稳定点，以稳定点构建动态参考并执行停车确认、占用维持与释放判定。

\subsection{稳定窗判定与稳定点提取}
\label{subsec:ch4_stable}

停车与占用语义的关键在于稳态水平而非瞬时差分。本文将稳定点定义为稳定窗均值向量，并采用窗内波动与窗间水平差联合约束，以抑制缓行或停走引起的伪稳态误判。

令稳定窗长度为$L$，以采样点$k$结束的滑动窗口记为$\mathcal{W}(k)=\{k-L+1,\ldots,k\}$，其窗口均值向量为
\begin{equation}
\overline{\mathbf{B}}(k)=\frac{1}{L}\sum_{j\in\mathcal{W}(k)}\widetilde{\mathbf{B}}(j).
\label{eq:ch4_mean}
\end{equation}

定义三轴窗内极差向量$\mathbf{R}(k)$与窗间均值差向量$\mathbf{M}(k)$为
\begin{equation}
\mathbf{R}(k)=
\begin{bmatrix}
\max\limits_{j \in \mathcal{W}(k)}\widetilde{B}_x(j) - \min\limits_{j \in \mathcal{W}(k)}\widetilde{B}_x(j)\\
\max\limits_{j \in \mathcal{W}(k)}\widetilde{B}_y(j) - \min\limits_{j \in \mathcal{W}(k)}\widetilde{B}_y(j)\\
\max\limits_{j \in \mathcal{W}(k)}\widetilde{B}_z(j) - \min\limits_{j \in \mathcal{W}(k)}\widetilde{B}_z(j)
\end{bmatrix},
\qquad
\mathbf{M}(k)=\overline{\mathbf{B}}(k)-\overline{\mathbf{B}}(k-s),
\label{eq:ch4_range_mean_vec}
\end{equation}
其中步长$s\ge 1$，工程实现中默认取$s=L$以比较相邻时间段的平均水平。

为减弱轴间尺度差异对阈值设置的影响，引入权重向量$\mathbf{w}=[w_x,w_y,w_z]^{\mathsf{T}}$，并对极差与均值差进行加权融合：
\begin{equation}
R(k)=\mathbf{w}^{\mathsf{T}}\mathbf{R}(k),\qquad
M(k)=\mathbf{w}^{\mathsf{T}}\bigl|\mathbf{M}(k)\bigr|,
\label{eq:ch4_RM_weight}
\end{equation}
其中$|\cdot|$表示逐元素取绝对值。权重$\mathbf{w}$可根据无车稳定段三轴磁场的标准差$\sigma_x,\sigma_y,\sigma_z$设定，
并归一化使$w_x+w_y+w_z=1$，以抑制噪声较大的轴对融合量的主导效应。

稳定窗判据定义为
\begin{equation}
R(k)\le R_{\mathrm{th}},\qquad M(k)\le M_{\mathrm{th}}.
\label{eq:ch4_stable_cond}
\end{equation}

为抑制阈值附近抖动，采用连续门控实现稳定判定。定义稳定计数器$n_{\mathrm{st}}(k)$为
\begin{equation}
n_{\mathrm{st}}(k)=
\begin{cases}
\min\{n_{\mathrm{st}}(k-1)+1,\,N_{\mathrm{stable}}\}, & R(k)\le R_{\mathrm{th}}\ \land\ M(k)\le M_{\mathrm{th}},\\
0, & \text{其他情况},
\end{cases}
\label{eq:ch4_nst}
\end{equation}
并定义稳定状态$g_{\mathrm{st}}(k)$与稳定点触发标志$f_{\mathrm{st}}(k)$为
\begin{equation}
g_{\mathrm{st}}(k)=\mathbb{I}\{n_{\mathrm{st}}(k)=N_{\mathrm{stable}}\},\qquad
f_{\mathrm{st}}(k)=g_{\mathrm{st}}(k)\cdot \bigl(1-g_{\mathrm{st}}(k-1)\bigr),
\label{eq:ch4_stable_flag}
\end{equation}
其中$\mathbb{I}\{\cdot\}$为指示函数。$f_{\mathrm{st}}(k)=1$表示稳定状态的起始时刻，用于锁定稳定点：
\begin{equation}
\mathbf{S}(k)=\overline{\mathbf{B}}(k),\qquad \text{当 } f_{\mathrm{st}}(k)=1.
\label{eq:ch4_S_lock}
\end{equation}
工程实现中，每段稳定区间仅锁定一次稳定点以避免重复触发。

\subsection{动态参考构建与漂移及相似性判定}
\label{subsec:ch4_ref}

稳定点用于构建动态参考，以表达不同语义状态下的稳态水平。本文维护两类动态参考：环境参考$\mathbf{S}_{\mathrm{pre}}$与占用参考$\mathbf{S}_{\mathrm{post}}$。

\paragraph{环境参考的更新}
当系统处于空闲状态（\texttt{occ\_flag}=0）且锁定稳定点$\mathbf{S}(k)$时，更新环境参考为
\begin{equation}
\mathbf{S}_{\mathrm{pre}}\leftarrow \mathbf{S}(k).
\label{eq:ch4_pre_update}
\end{equation}
该更新仅在无车稳定段执行，用以避免车辆扰动对环境参考的污染。

\paragraph{停车确认与占用参考锁定}
对第$m$次事件$\mathcal{E}_m$，当到达其结束时刻$t_{\mathrm{out},m}$后，在寻稳时长$T_{\mathrm{seek}}$内搜索稳定点。
若在区间$[t_{\mathrm{out},m},\,t_{\mathrm{out},m}+T_{\mathrm{seek}}]$内获得停车后稳定点，则记为$\mathbf{S}_{\mathrm{post}}$并定义漂移为
\begin{equation}
\Delta\mathbf{B}=\mathbf{S}_{\mathrm{post}}-\mathbf{S}_{\mathrm{pre}},\qquad
\Delta B=\lVert\Delta\mathbf{B}\rVert_2.
\label{eq:ch4_drift}
\end{equation}
当$\Delta B>D_{\mathrm{th}}$时确认停车成立，置\texttt{park\_flag}=1并进入占用状态（\texttt{occ\_flag}=1），
并将$t_{\mathrm{park,in}}$对齐至$t_{\mathrm{out},m}$。

\paragraph{占用维持、释放判定与参考维护}
在\texttt{OCCUPIED}阶段，通行车辆可能引发新的事件片段$\mathcal{E}_n$。当事件结束于$t_{\mathrm{out},n}$后，
若获得新稳定点$\mathbf{S}_{\mathrm{new}}$，定义
\begin{equation}
\Delta\mathbf{B}_{\mathrm{new}}=\mathbf{S}_{\mathrm{new}}-\mathbf{S}_{\mathrm{pre}},\qquad
\mathrm{dist}=\left\lVert \Delta\mathbf{B}_{\mathrm{new}}-\Delta\mathbf{B}\right\rVert_2.
\label{eq:ch4_dist}
\end{equation}

释放判定采用回归环境参考优先的规则：
\begin{enumerate}
  \item 若$\lVert\mathbf{S}_{\mathrm{new}}-\mathbf{S}_{\mathrm{pre}}\rVert_2<D_{\mathrm{free}}$，判定释放成立，
  置\texttt{occ\_flag}=0并记录$t_{\mathrm{park,out}}=t_{\mathrm{out},n}$；
  \item 否则若$\mathrm{dist}<\mathrm{dist}_{\mathrm{th}}$，判定仍处于占用稳态（多为过车干扰后回归占用稳态）；
  \item 否则转入候选释放状态，由外层状态机结合超时机制与后续证据完成判决。
\end{enumerate}

为提高长时占用下对慢漂移的适应性，当判定仍占用且满足$\mathrm{dist}<\mathrm{dist}_{\mathrm{th}}$时，可对占用参考进行更新：
\begin{equation}
\mathbf{S}_{\mathrm{post}}\leftarrow (1-\lambda)\mathbf{S}_{\mathrm{post}}+\lambda\,\mathbf{S}_{\mathrm{new}},\qquad 0<\lambda\le 1.
\label{eq:ch4_post_ema}
\end{equation}
同时，为补偿环境慢漂移对释放阈值的影响，可选地将占用参考的变化量修正到环境参考中：
\begin{equation}
\mathbf{S}_{\mathrm{pre}}\leftarrow \mathbf{S}_{\mathrm{pre}}+\bigl(\mathbf{S}_{\mathrm{post}}-\mathbf{S}_{\mathrm{post,old}}\bigr),
\label{eq:ch4_pre_comp}
\end{equation}
其中$\mathbf{S}_{\mathrm{post,old}}$为更新前的占用参考。该策略的前提是变化由环境漂移主导，因而需在鲁棒性设计中对更新幅度与触发条件设置门控。

\subsection{外层状态机与连续车流退化分支}
\label{subsec:ch4_fsm}

为将停车与占用语义显式化并与第2章事件触发形成可复现的判决链条，本文在稳定点与判据之上构建外层有限状态机。
状态集合包括\texttt{FREE}、\texttt{CANDIDATE\_PARK}、\texttt{OCCUPIED}与\texttt{CANDIDATE\_RELEASE}。
其中，\texttt{CANDIDATE\_PARK}用于表达事件结束后等待稳态证据的阶段，\texttt{CANDIDATE\_RELEASE}用于表达占用态下出现不一致证据后的待确认阶段。
状态转移由寻稳超时$T_{\mathrm{seek}}$、阈值滞回以及一致性计数器共同抑制临界抖动。

连续车流或强扰动背景下，事件结束后可能在$T_{\mathrm{seek}}$内无法获得稳定点。为避免确认延迟过大或漏检，
本文在寻稳超时后启用退化处理：以固定窗均值近似稳态，并以多次一致性证据替代单次稳定点证据。

设事件结束采样点索引为$k_{\mathrm{out}}$，搜索窗口末端索引为
\begin{equation}
k_{\mathrm{seek}} = k_{\mathrm{out}} + \lfloor T_{\mathrm{seek}} f_s \rfloor.
\label{eq:ch4_kseek}
\end{equation}
设固定窗长度为$L_{\mathrm{fix}}$，取末端固定窗均值作为近似稳态
\begin{equation}
\hat{\mathbf{S}} = \frac{1}{L_{\mathrm{fix}}}\sum_{j = k_{\mathrm{seek}} - L_{\mathrm{fix}} + 1}^{k_{\mathrm{seek}}}\widetilde{\mathbf{B}}(j).
\label{eq:ch4_Shat}
\end{equation}

\paragraph{候选停车阶段的退化确认}
当\texttt{CANDIDATE\_PARK}中稳定点不可得时，计算近似漂移$\Delta\hat{\mathbf{B}}=\hat{\mathbf{S}}-\mathbf{S}_{\mathrm{pre}}$及幅值$\Delta\hat{B}=\lVert\Delta\hat{\mathbf{B}}\rVert_2$。
定义一致性计数器递推为
\begin{equation}
c_n \leftarrow
\begin{cases}
\min\{c_n+1,\,c_{\mathrm{th}}\}, & \Delta\hat{B} > D_{\mathrm{th}},\\
0, & \text{其他情况},
\end{cases}
\label{eq:ch4_cnt_park}
\end{equation}
当$c_n=c_{\mathrm{th}}$时确认停车成立，并进入\texttt{OCCUPIED}状态。若在超时窗口内始终无法累积到门限，则判定该次事件为通过或干扰，不进入占用状态。

\paragraph{候选释放阶段的退化确认}
当\texttt{CANDIDATE\_RELEASE}中稳定点不可得时，采用回归环境参考作为释放证据。定义一致性计数器递推为
\begin{equation}
c_n \leftarrow
\begin{cases}
\min\{c_n+1,\,c_{\mathrm{th}}\}, & \lVert\hat{\mathbf{S}}-\mathbf{S}_{\mathrm{pre}}\rVert_2 < D_{\mathrm{free}},\\
0, & \text{其他情况},
\end{cases}
\label{eq:ch4_cnt_free}
\end{equation}
当$c_n=c_{\mathrm{th}}$时确认释放成立并回到\texttt{FREE}状态。若后续获得稳定点且相似性满足占用条件，则计数器清零并回退到\texttt{OCCUPIED}。

上述退化分支通过$L_{\mathrm{fix}}$与$c_{\mathrm{th}}$的组合抑制瞬时窗污染导致的误判，其中$c_{\mathrm{th}}>1$可有效降低固定窗被单次过车干扰污染时的误触发风险。

\section{鲁棒性设计与参数设置}
\label{sec:ch4_robust}

上一节给出了稳定点提取、动态参考构建、漂移与相似性判据以及外层状态机逻辑。本节从工程鲁棒性角度进一步讨论关键退化情形与对应策略。道路侧单地磁场景中，鲁棒性风险主要来自三类因素：其一，环境慢漂移与偶发阶跃干扰使参考维护面临误更新风险；其二，缓行、停走或邻道弱扰动可能造成“变化变慢”的伪稳态，从而触发错误的稳定点锁定；其三，连续车流背景下稳定窗缺失，使停车确认与释放判定缺乏一次性稳态证据。围绕上述问题，本文在稳定点判据、参考更新门控以及退化分支设计上引入多重约束，并给出参数设置的可复现实证路径。

\subsection{环境漂移抑制与参考更新门控}
\label{subsec:ch4_drift_gate}

环境磁场在长时间运行中会发生缓慢漂移，此外道路附近的外源电磁环境亦可能带来阶跃型扰动。若环境参考更新条件过宽，容易在非空车稳态或异常稳态下发生误更新，进而诱发误判或误释放。为此，本文对环境参考$\mathbf{S}_{\mathrm{pre}}$采用严格门控更新：仅当外层状态处于\texttt{FREE}，且稳定判据满足并通过更新幅度安全门控时，才允许更新。

令$\mathbf{S}(k)$为稳定点候选（上一节中由稳定段锁定得到），定义更新门控指示量
\begin{equation}
g_{\mathrm{upd}}(k) =
\mathbb{I}\!\left(\texttt{state}=\texttt{FREE}\right)\cdot
\mathbb{I}\!\left(f_{\mathrm{st}}(k)=1\right)\cdot
\mathbb{I}\!\left(\left\lVert \mathbf{S}(k)-\mathbf{S}_{\mathrm{pre}}\right\rVert_2 < D_{\mathrm{upd}}\right),
\label{eq:ch4_upd_gate}
\end{equation}
其中$D_{\mathrm{upd}}$为单次更新幅度上限，用于抑制阶跃干扰导致的误更新。当$g_{\mathrm{upd}}(k)=1$时，环境参考采用指数滑动平均更新：
\begin{equation}
\mathbf{S}_{\mathrm{pre}}^{+} = (1-\alpha)\mathbf{S}_{\mathrm{pre}}^{-} + \alpha\,\mathbf{S}(k),\qquad 0<\alpha<1.
\label{eq:ch4_ema_update_refine}
\end{equation}
其中$\alpha$控制慢漂移跟踪速度，$\alpha$越大，参考跟踪越快，但对短时波动更敏感；$\alpha$越小，更新更保守但对漂移响应更慢。

占用参考$\mathbf{S}_{\mathrm{post}}$在上一节中默认于停车确认后冻结，以维持占用稳态水平的可解释性。考虑到长时占用过程中环境仍可能发生慢漂移，为避免占用参考与环境参考产生系统性偏移，本文提供一种可选的“占用态漂移补偿”策略：当处于\texttt{OCCUPIED}且事件后获得新的稳定点$\mathbf{S}_{\mathrm{new}}$，并通过相似性判据确认仍为同一占用稳态时，对$\mathbf{S}_{\mathrm{post}}$进行保守更新，同时将其增量同步补偿到$\mathbf{S}_{\mathrm{pre}}$，以抵消环境慢漂移的累积影响。具体地，若$\mathrm{dist}<\mathrm{dist}_{\mathrm{th}}$且
\begin{equation}
\left\lVert \mathbf{S}_{\mathrm{new}}-\mathbf{S}_{\mathrm{post}}\right\rVert_2 < D_{\mathrm{comp}},
\label{eq:ch4_comp_gate}
\end{equation}
则执行
\begin{align}
\mathbf{S}_{\mathrm{post}}^{+} &= (1-\alpha_{\mathrm{occ}})\mathbf{S}_{\mathrm{post}}^{-} + \alpha_{\mathrm{occ}}\,\mathbf{S}_{\mathrm{new}}, \label{eq:ch4_post_update}\\
\mathbf{S}_{\mathrm{pre}}^{+}  &= \mathbf{S}_{\mathrm{pre}}^{-} + \gamma\left(\mathbf{S}_{\mathrm{post}}^{+}-\mathbf{S}_{\mathrm{post}}^{-}\right), \qquad 0\le \gamma \le 1,
\label{eq:ch4_pre_comp}
\end{align}
其中$\alpha_{\mathrm{occ}}$通常取小于$\alpha$的保守值，$D_{\mathrm{comp}}$限制占用态更新的最大幅度，$\gamma$用于控制补偿强度（工程上可取$\gamma=1$以实现等量补偿）。该策略的目的在于：在不破坏占用稳态语义的前提下，使相似性判据对慢漂移更不敏感，从而降低长时占用下的误释放概率。

\subsection{伪稳态与误触发抑制：双判据与时序约束}
\label{subsec:ch4_pseudo}

在缓行、停走或车流密度较大场景中，信号变化速率可能显著降低，窗内极差$R(k)$可能变小，但窗口均值仍持续缓慢漂移。此类“变化变慢”并不等价于“进入稳态”，若仅用窗内波动判据，容易将缓慢变化误判为稳定点并导致参考误更新或误确认。因此，本文采用窗内波动$R(k)$与窗间水平差$M(k)$联合约束（见式\eqref{eq:ch4_stable_cond}），其含义分别为：$R(k)$刻画窗口内的幅值波动范围，$M(k)$刻画当前窗口与历史窗口的水平差异，两者共同用于区分缓行引起的缓慢变化与环境慢漂移或真正稳态。

在联合判据基础上，采用连续门控计数$n_{\mathrm{st}}(k)$实现最短持续时间约束（式\eqref{eq:ch4_nst}），从时序上抑制阈值附近抖动导致的频繁进出稳态。为进一步提升工程稳定性，可在实现中引入阈值滞回：使用进入阈值$(R_{\mathrm{th}}^{\mathrm{in}},M_{\mathrm{th}}^{\mathrm{in}})$锁定稳定段，使用更宽松的退出阈值$(R_{\mathrm{th}}^{\mathrm{out}},M_{\mathrm{th}}^{\mathrm{out}})$解除稳定段，其中$R_{\mathrm{th}}^{\mathrm{out}}>R_{\mathrm{th}}^{\mathrm{in}}$且$M_{\mathrm{th}}^{\mathrm{out}}>M_{\mathrm{th}}^{\mathrm{in}}$，以减少稳态边界处的重复触发。

此外，释放判定中引入$D_{\mathrm{free}}<D_{\mathrm{th}}$的滞回约束：停车确认要求占用漂移幅值超过$D_{\mathrm{th}}$，而释放要求回归环境参考小于$D_{\mathrm{free}}$，从而降低阈值附近抖动与弱扰动导致的误释放风险。

\subsection{连续车流工况：退化分支的触发与参数权衡}
\label{subsec:ch4_degrade_param}

连续车流或强扰动背景下，事件结束后可能在$T_{\mathrm{seek}}$时间内无法获得满足稳定判据的窗口，从而$\mathbf{S}_{\mathrm{post}}$与$\mathbf{S}_{\mathrm{new}}$不可得。若完全依赖稳定点证据，停车确认与释放判定会出现确认延迟增大甚至漏检。为维持算法可用性，本文在外层状态机中设置寻稳超时$T_{\mathrm{seek}}$：当超时发生时启用退化分支，以固定窗均值近似稳态，并以一致性计数器将单次稳态证据转化为多次事件的一致证据。

设事件结束采样点为$k_{\mathrm{out}}$，搜索窗口末端为$k_{\mathrm{seek}}$（式\eqref{eq:ch4_kseek}）。固定窗近似稳态定义为$\hat{\mathbf{S}}$（式\eqref{eq:ch4_Shat}）。在\texttt{CANDIDATE\_PARK}或\texttt{OCCUPIED}等需要稳态证据的阶段，构造近似相似性度量：
\begin{equation}
\widehat{\mathrm{dist}} = \left\lVert \left(\hat{\mathbf{S}}-\mathbf{S}_{\mathrm{pre}}\right) - \Delta\mathbf{B}\right\rVert_2
= \left\lVert \hat{\mathbf{S}}-\mathbf{S}_{\mathrm{post}}\right\rVert_2 .
\label{eq:ch4_hatdist}
\end{equation}
当$\widehat{\mathrm{dist}}<\mathrm{dist}_{\mathrm{th}}$时认为该次事件结束后的近似稳态与占用稳态一致，计数器增加；否则计数器清零。记一致性计数器为$c_n$，其更新可写为
\begin{equation}
c_n \leftarrow
\begin{cases}
\min\{c_n+1,\,c_{\mathrm{th}}\}, & \widehat{\mathrm{dist}}<\mathrm{dist}_{\mathrm{th}},\\
0, & \text{其他情况}.
\end{cases}
\label{eq:ch4_cn_update}
\end{equation}
当$c_n=c_{\mathrm{th}}$时，确认停车或确认仍占用。释放退化判定同理：将一致性条件替换为回归条件$\lVert \hat{\mathbf{S}}-\mathbf{S}_{\mathrm{pre}}\rVert_2 < D_{\mathrm{free}}$，达到计数门限后确认释放。

参数权衡方面：$T_{\mathrm{seek}}$越大越可能获得真实稳定点，但确认延迟增大；$L_{\mathrm{fix}}$越大平滑更强但滞后更大；$c_{\mathrm{th}}$越大越稳健但确认更慢。本文在实验部分分别统计拥堵组数据的漏检率、误报率与确认延迟，以综合指标进行折中选择。

\subsection{关键参数确定方法（概述）}
\label{subsec:ch4_paramset}

本文参数采用数据驱动方式确定，并遵循“先稳态判据，再漂移判据，最后退化分支”的顺序，以降低耦合调参难度。

\textbf{1）稳定窗相关参数$L,s,R_{\mathrm{th}},M_{\mathrm{th}},N_{\mathrm{stable}}$：}
首先在\texttt{FREE}状态下选取无车稳态数据段，统计$R(k)$与$M(k)$的经验分布，并将$R_{\mathrm{th}}$与$M_{\mathrm{th}}$取为高分位数以控制稳态误触发概率。随后在缓行、停走等伪稳态数据段统计$M(k)$分布，并调小$M_{\mathrm{th}}$或增大$N_{\mathrm{stable}}$，以抑制慢变伪稳态。窗口步长$s$通常取$s=L$以比较相邻窗口水平差，窗口长度$L$依据采样率与期望最短稳态持续时间选取。

\textbf{2）停车确认阈值$D_{\mathrm{th}}$与释放阈值$D_{\mathrm{free}}$：}
在包含“正常通过”与“真实停车”的标注数据上统计漂移幅值$\Delta B$分布，对$D_{\mathrm{th}}$进行阈值扫描，优先满足“通过事件不被误判为停车”的低误报约束，并在此基础上选择验证集综合指标最优的阈值。释放阈值$D_{\mathrm{free}}$基于无车稳态噪声水平确定，并保持$D_{\mathrm{free}}<D_{\mathrm{th}}$形成滞回。

\textbf{3）相似性阈值$\mathrm{dist}_{\mathrm{th}}$：}
在\texttt{OCCUPIED}阶段分别收集“占用态过车后回到占用稳态”的样本与“真实释放后回到环境稳态”的样本，统计$\mathrm{dist}$或$\lVert\mathbf{S}_{\mathrm{new}}-\mathbf{S}_{\mathrm{post}}\rVert_2$的分布差异，通过分布分离度或ROC曲线确定$\mathrm{dist}_{\mathrm{th}}$。

\textbf{4）参考更新参数$\alpha,D_{\mathrm{upd}}$与占用态补偿参数$\alpha_{\mathrm{occ}},D_{\mathrm{comp}},\gamma$：}
$\alpha$用于环境慢漂移跟踪，其取值需在“漂移可跟踪”与“短时扰动不过度跟随”之间折中；$D_{\mathrm{upd}}$用于抑制阶跃型异常稳态导致的误更新，宜显著小于典型车辆占用漂移幅值。若启用占用态漂移补偿，建议取较小$\alpha_{\mathrm{occ}}$并设置$D_{\mathrm{comp}}$上限，同时令$\gamma$不大于1，以降低对偶发扰动的敏感性。

\textbf{5）退化分支参数$T_{\mathrm{seek}},L_{\mathrm{fix}},c_{\mathrm{th}}$：}
在拥堵与连续车流数据上联合调参，以漏检率与确认延迟为约束进行折中。一般而言，$T_{\mathrm{seek}}$应覆盖若干个典型车间隙以提高获得稳定窗的概率；当稳定窗长期不可得时，$L_{\mathrm{fix}}$用于保证固定窗均值的估计方差可控，$c_{\mathrm{th}}$用于将判据从单次证据提升为多次一致证据，从而提高连续车流背景下的鲁棒性。


\section{实验设置与结果分析}
\label{sec:ch4_exp}

本节基于道路侧单地磁三轴节点的实测数据，对第\ref{sec:ch4_method}节提出的停车与占用判定方法进行验证。实验关注三类指标：其一，占用状态段的检测性能（精确率、召回率与$F_1$）；其二，占用起止边界与占用时长的误差统计；其三，停车确认与释放确认的上报延迟，以及关键模块在复杂工况下的消融贡献。

\subsection{数据采集、标注流程与工况分组}
\label{subsec:ch4_data}

\subsubsection{测试场景与部署方式}

实验数据由道路侧单地磁三轴传感节点采集，采样率为$f_s=50\,\mathrm{Hz}$。节点敏感面与路面齐平安装，安装位置位于路肩或应急车道侧的近车道线区域，使其对目标车道停靠车辆具有足够敏感度，同时能够在主车道通行背景下采集到过车干扰。

为覆盖“停靠与连续通行并存”的典型道路工况，测试场景优先选择双车道或含路肩的路段：右侧车道（或路肩、应急车道）用于产生停靠与占用事件，邻近车道保持车辆连续通行，用于模拟占用态下的过车扰动与拥堵背景。测试期间通过视频设备同步记录现场情况，用于真值复核与标注对齐。

\subsubsection{坐标定义与数据对齐}

建立节点坐标系：$x$轴指向车辆行驶方向，$y$轴指向道路外侧，$z$轴竖直向上。采集端为每个采样点记录时间戳，并对视频帧时间与采样时间进行对齐。预处理输出与第2章保持一致，即对每个采样点得到预处理后三轴序列$\widetilde{\mathbf{B}}(k)$，并由第2章事件检测底座输出\texttt{diff\_on}/\texttt{diff\_off}对应的$(t_{\mathrm{in}},t_{\mathrm{out}})$。

对每次真实停靠事件，标注真值占用区间
\begin{equation}
[t^{*}_{\mathrm{park,in}},\,t^{*}_{\mathrm{park,out}}],
\end{equation}
其中$t^{*}_{\mathrm{park,in}}$定义为车辆完成停稳并形成持续占用的起始时刻，$t^{*}_{\mathrm{park,out}}$定义为车辆离开敏感范围且信号回归环境稳态的时刻。上述定义强调“形成稳态占用”与“回归环境稳态”，与本章稳定点判据的语义口径一致。

\subsubsection{工况分组}

为覆盖道路停车检测中影响判定可靠性的关键因素，本文将数据划分为四类工况：
\begin{itemize}
  \item A类：正常车流下的单车停靠，占用段前后可获得稳定窗；
  \item B类：占用存在且伴随过车干扰，占用段内出现若干事件片段，但事件后能够回到占用稳态；
  \item C类：连续车流或拥堵导致稳定窗缺失，事件结束后在$T_{\mathrm{seek}}$内难以获得稳定点；
  \item D类：存在环境慢漂移或阶跃型强干扰，容易触发参考误更新与误释放。
\end{itemize}

各组数据统计如表\ref{tab:ch4_dataset}所示。为避免阈值过拟合，本文将数据划分为验证集与测试集：验证集用于阈值与参数选择，测试集用于最终报告，测试阶段所有方法参数固定不再调整。

\begin{table}[t]
\centering
\caption{数据集与工况分组统计}
\label{tab:ch4_dataset}
\begin{tabular}{c r r r r p{0.32\linewidth}}
\hline
组别 & 数据段数 & 总时长(min) & 占用事件数 & 通过事件数 & 说明 \\
\hline
A &  &  &  &  & 正常车流单车停靠，稳定窗可得 \\
B &  &  &  &  & 占用存在且伴随过车干扰 \\
C &  &  &  &  & 连续车流或拥堵（稳定窗缺失） \\
D &  &  &  &  & 漂移或强干扰背景 \\
\hline
\end{tabular}
\end{table}

\subsection{对比方法与评价指标}
\label{subsec:ch4_metric}

\subsubsection{对比方法}

为验证本文方法的有效性与鲁棒性，设置两条代表性基线进行对比：
\begin{enumerate}
  \item Baseline-1（单参考与幅值滞回）：仅维护环境参考（\texttt{FREE}状态更新），在事件结束后以固定窗均值与环境参考的幅值差进行停车确认与释放判定，不显式构造占用参考与相似性判据；
  \item Baseline-2（事件驱动固定窗近似）：在\texttt{diff\_off}后等待固定时窗并取均值近似稳态，完成停车确认与释放判定，但不包含稳定窗双判据、占用漂移相似性度量与连续车流一致性退化机制；
  \item Ours：本文提出的稳定点提取、动态参考、漂移与相似性判据，以及包含连续车流退化分支的外层状态机。
\end{enumerate}
三种方法使用相同的预处理输出与相同的数据划分方式；阈值均在验证集上选择，并在测试集上固定。

\subsubsection{区间匹配与检测性能指标}

以占用状态段为评价对象。设真值占用区间为$[t^{*}_{\mathrm{park,in}},t^{*}_{\mathrm{park,out}}]$，算法输出占用区间为$[t_{\mathrm{park,in}},t_{\mathrm{park,out}}]$，定义区间交并比
\begin{equation}
\mathrm{IoU}=
\frac{\mathrm{len}\!\left([t_{\mathrm{park,in}}, t_{\mathrm{park,out}}]\cap[t^{*}_{\mathrm{park,in}}, t^{*}_{\mathrm{park,out}}]\right)}
{\mathrm{len}\!\left([t_{\mathrm{park,in}}, t_{\mathrm{park,out}}]\cup[t^{*}_{\mathrm{park,in}}, t^{*}_{\mathrm{park,out}}]\right)}.
\end{equation}
采用一对一匹配策略：对每个预测区间，寻找与其$\mathrm{IoU}$最大的真值区间进行候选匹配，并按$\mathrm{IoU}$从大到小贪心锁定匹配对；当$\mathrm{IoU}\ge\theta$时记为TP，否则记为FP；未被匹配的真值区间记为FN。本文默认$\theta=0.5$。据此定义
\begin{equation}
\mathrm{Precision}=\frac{TP}{TP+FP},\qquad
\mathrm{Recall}=\frac{TP}{TP+FN},\qquad
F_1=\frac{2\,\mathrm{Precision}\cdot\mathrm{Recall}}{\mathrm{Precision}+\mathrm{Recall}}.
\end{equation}

为便于工程口径对照，本文同时报告误检率与漏检率：
\begin{equation}
R_{\mathrm{f}}=\frac{FP}{N_{\mathrm{gt}}}\times 100\%,\qquad
R_{\mathrm{m}}=\frac{FN}{N_{\mathrm{gt}}}\times 100\%,
\end{equation}
其中$N_{\mathrm{gt}}$为真值占用事件总数。该口径与停车检测相关工作中常用的“误检次数、漏检次数及其比例”一致，便于直接反映误报与漏报水平。

\subsubsection{边界误差、时长误差与上报延迟}

对TP集合统计边界误差与占用时长误差：
\begin{equation}
\Delta t_{\mathrm{in}} = t_{\mathrm{park,in}} - t^{*}_{\mathrm{park,in}},\quad
\Delta t_{\mathrm{out}} = t_{\mathrm{park,out}} - t^{*}_{\mathrm{park,out}},\quad
\Delta T = (t_{\mathrm{park,out}} - t_{\mathrm{park,in}}) - (t^{*}_{\mathrm{park,out}} - t^{*}_{\mathrm{park,in}}).
\end{equation}
由于本文将$t_{\mathrm{park,in}}$与$t_{\mathrm{park,out}}$对齐至相应扰动结束时刻$t_{\mathrm{out}}$作为边界近似，算法的实际\emph{上报}时刻可能晚于边界估计时刻。为评估“尽可能早上报”的能力，额外定义停车确认与释放确认的上报延迟：
\begin{equation}
\tau_{\mathrm{in}} = t_{\mathrm{conf,in}} - t^{*}_{\mathrm{park,in}},\qquad
\tau_{\mathrm{out}} = t_{\mathrm{conf,out}} - t^{*}_{\mathrm{park,out}},
\end{equation}
其中$t_{\mathrm{conf,in}}$为\texttt{park\_flag}触发时刻，$t_{\mathrm{conf,out}}$为\texttt{occ\_flag}由1变为0的确认时刻。本文报告上述误差与延迟的均值、中位数及$95\%$分位数。

\subsection{对比结果、消融与讨论}
\label{subsec:ch4_results}

\subsubsection{总体与分工况性能}

表\ref{tab:ch4_bycase}给出分工况与总体性能对比。图\ref{fig:ch4_cases}展示四类工况的典型样例，包括三轴信号、稳定标志$f_{\mathrm{st}}(k)$、稳定点序列、状态机状态与真值区间叠加，并突出连续车流退化分支的触发与一致性确认过程。

\begin{table}[t]
\centering
\caption{分工况与总体性能对比（$\theta=0.5$）}
\label{tab:ch4_bycase}
\begin{tabular}{c l c c c c c p{0.26\linewidth}}
\hline
工况 & 方法 & Precision & Recall & $F_1$ & $R_{\mathrm{f}}(\%)$ & $R_{\mathrm{m}}(\%)$ & 主要错误形态 \\
\hline
A & Baseline-1 &  &  &  &  &  &  \\
A & Baseline-2 &  &  &  &  &  &  \\
A & Ours       &  &  &  &  &  &  \\
\hline
B & Baseline-1 &  &  &  &  &  &  \\
B & Baseline-2 &  &  &  &  &  &  \\
B & Ours       &  &  &  &  &  &  \\
\hline
C & Baseline-1 &  &  &  &  &  &  \\
C & Baseline-2 &  &  &  &  &  &  \\
C & Ours       &  &  &  &  &  &  \\
\hline
D & Baseline-1 &  &  &  &  &  &  \\
D & Baseline-2 &  &  &  &  &  &  \\
D & Ours       &  &  &  &  &  &  \\
\hline
All & Baseline-1 &  &  &  &  &  &  \\
All & Baseline-2 &  &  &  &  &  &  \\
All & Ours       &  &  &  &  &  &  \\
\hline
\end{tabular}
\end{table}

% \begin{figure}[t]
% \centering
% \begin{tabular}{p{0.48\linewidth} p{0.48\linewidth}}
% \centering \includegraphics[width=0.95\linewidth]{figures/ch4_case_A.pdf} &
% \centering \includegraphics[width=0.95\linewidth]{figures/ch4_case_B.pdf} \\
% (a) A类：正常车流单车停靠 &
% (b) B类：占用伴随过车干扰 \\
% \centering \includegraphics[width=0.95\linewidth]{figures/ch4_case_C.pdf} &
% \centering \includegraphics[width=0.95\linewidth]{figures/ch4_case_D.pdf} \\
% (c) C类：连续车流（稳定窗缺失） &
% (d) D类：漂移或强干扰背景 \\
% \end{tabular}
% \caption{四类工况典型样例与算法输出叠加。}
% \label{fig:ch4_cases}
% \end{figure}

\subsubsection{边界误差与上报延迟分析}

表\ref{tab:ch4_timing}统计边界误差与上报延迟。由于本文边界估计对齐至扰动结束时刻$t_{\mathrm{out}}$，因此$\Delta t_{\mathrm{in}}$通常呈现正偏；而上报延迟$\tau_{\mathrm{in}}$还包含寻稳与门控计数带来的额外等待。对比各方法可验证稳定窗双判据、相似性度量以及退化分支对延迟与误差的影响。

\begin{table}[t]
\centering
\caption{边界误差与上报延迟统计（TP集合）}
\label{tab:ch4_timing}
\begin{tabular}{l c c c c c c}
\hline
方法 & $\mathrm{median}(|\Delta t_{\mathrm{in}}|)$ & $\mathrm{median}(|\Delta t_{\mathrm{out}}|)$ &
$\mathrm{median}(|\Delta T|)$ &
$\mathrm{median}(\tau_{\mathrm{in}})$ & $\mathrm{median}(\tau_{\mathrm{out}})$ &
$P_{95}(\tau_{\mathrm{in}})$ \\
\hline
Baseline-1 &  &  &  &  &  &  \\
Baseline-2 &  &  &  &  &  &  \\
Ours       &  &  &  &  &  &  \\
\hline
\end{tabular}
\end{table}

\subsubsection{消融实验与讨论}

为验证关键模块对复杂工况的贡献，设置如下消融变体：
\begin{itemize}
  \item Abl-1：去除均值差判据$M(k)$，仅保留窗内波动$R(k)$；
  \item Abl-2：去除相似性度量$\mathrm{dist}$，占用态仅使用回归阈值进行释放判定；
  \item Abl-3：去除连续车流退化分支（固定窗近似与一致性计数器）；
  \item Abl-4：去除参考更新安全门（式\eqref{eq:ch4_safe_update}）。
\end{itemize}
消融结果见表\ref{tab:ch4_ablation}，并结合C类与D类样例对误检与漏检原因进行解释：去除$M(k)$会显著增加缓行伪稳态触发；去除$\mathrm{dist}$会在占用态过车干扰后产生误释放或重复占用；去除退化分支会在拥堵场景下引起确认延迟显著增大并导致漏检；去除更新安全门会在阶跃干扰下产生参考误更新，进而引发连锁误判。

\begin{table}[t]
\centering
\caption{消融实验结果（$F_1$）}
\label{tab:ch4_ablation}
\begin{tabular}{l c c c c p{0.30\linewidth}}
\hline
变体 & $F_1$(All) & $F_1$(A) & $F_1$(B) & $F_1$(C) & 主要退化现象 \\
\hline
Ours-full &  &  &  &  &  \\
Abl-1（去除$M$） &  &  &  &  & 伪稳态误触发增多，误检上升 \\
Abl-2（去除$\mathrm{dist}$） &  &  &  &  & 占用态过车后误释放或重复占用 \\
Abl-3（去除退化分支） &  &  &  &  & 拥堵下寻稳失败导致漏检与延迟增大 \\
Abl-4（去除更新安全门） &  &  &  &  & 阶跃干扰下参考误更新，产生连锁误判 \\
\hline
\end{tabular}
\end{table}


\section{本章小结}
\label{sec:ch4_summary}

本章面向道路路肩或应急车道的异常停车与占用监测任务，在单地磁三轴信号（$50\,\mathrm{Hz}$）与端侧低算力约束下，研究如何在第2章车辆事件检测底座输出的事件片段之上，进一步形成可上报的停车确认、占用维持与释放边界等高层语义状态。针对停车停稳后差分量趋近于零、事件级输出难以直接表达占用持续态，以及连续车流背景下稳定窗缺失与伪稳态误触发等问题，本章提出了基于稳定点与动态参考的停车与占用判定框架。

在方法设计方面，本章主要完成了以下工作。第一，提出稳定窗判定与稳定点提取机制，将稳定点定义为稳定窗均值向量，并采用窗内波动与窗间水平差的联合约束及连续门控计数，抑制缓行、停走与弱扰动条件下的伪稳态误判。第二，构建稳定点驱动的动态参考体系，分别维护环境参考稳定点$\mathbf{S}_{\mathrm{pre}}$与占用参考稳定点$\mathbf{S}_{\mathrm{post}}$，并以稳定点漂移$\Delta\mathbf{B}$（及其幅值$\Delta B$）实现停车确认，以相似性度量$\mathrm{dist}$区分占用维持与释放，从而将停车与占用判定从瞬时事件提升为可持续更新的语义状态。第三，在上述判据之上构建外层有限状态机，引入必要的时序约束、阈值滞回与寻稳超时门控，保证状态转移链条可解释、可复现；针对连续车流导致稳定点不可得的退化情形，进一步引入固定窗近似与一致性计数器确认策略，使算法在拥堵与高车流背景下仍具备可用性与稳定性。

在实验验证方面，本章基于覆盖正常停靠、占用伴随过车干扰、连续车流稳定窗缺失以及漂移或强干扰等多类工况的数据进行评估，从占用状态段检测精度、边界误差与占用时长误差，以及停车确认与释放的时效性等维度对方法进行量化分析。对比实验验证了稳定点动态参考与相似性判据在占用维持与释放判定中的有效性；消融实验进一步表明，均值差门控对伪稳态抑制具有关键作用，连续车流退化分支能够显著提升拥堵工况下的可用性与鲁棒性，参考更新门控可降低漂移与异常干扰导致的参考误更新风险。

综上，本章通过稳定点、动态参考与外层状态机时序约束的组合，实现了停车与占用语义的端侧显式判定，形成了面向道路侧异常事件上报的可实施流程。上述方法与实验结论为后续章节的系统总结与进一步扩展（例如参数自适应、跨场景泛化与更复杂干扰条件下的鲁棒性提升）提供了基础。














































\chapter{基于眼视频运动补偿的三模态眼动事件检测}

\section{研究动机与多模态信息定义}

上一章的 \ModelMM{} 已经在 ``gaze + world-video'' 双模态设定下实现了较强的事件检测能力：gaze 分支通过希尔伯特空间编码与动力学特征注入建立对眼动轨迹的表达；world-video 分支通过凝视引导的高斯池化在 ``所看区域'' 提取局部外观与运动证据，并经时序差分聚合后与 gaze 在交叉注意力中融合，最终在眼动时间轴上输出逐点事件类别。该框架的关键前提是：\textbf{gaze 坐标既是运动学判别的主要信号，也是视频证据采样的空间锚点}。

在实际采集与跨被试 leave-one-out 泛化中，这个前提会受到系统误差的直接挑战。眼动事件检测高度依赖由差分构造的运动学信息：上一章中速度与方向等特征由 $(x_t,y_t)$ 的时间差分得到，差分对抖动噪声极其敏感，轻微的坐标抖动就可能诱发速度伪峰值，进而影响 GS 的边界定位与 GP 的持续段识别；与此同时，凝视引导高斯池化以 gaze 为采样中心，如果 gaze 存在系统性偏移或缓慢漂移，那么采样区域将长期偏离真实注视区域，导致 world-video 的证据被 ``读错''，再强的融合机制也无法弥补错误锚点带来的信息错配。也就是说，world-video 能帮助解释场景的外部变化，但它并不直接提供 ``如何修正 gaze 观测误差'' 的通道。

为缓解上述问题，本章引入第三模态 eye-video（单眼视频），并将其作用严格限定为 \textbf{gaze 运动补偿与矫正}：利用 eye-video 的时序变化学习一个幅度受限的补偿位移，并以门控方式注入到 gaze 中，从而在进入上一章的主干之前先提升 gaze 观测的稳定性。需要强调的是，本章数据没有瞳孔中心标注，且瞳孔近似处于 eye-video 的中心区域，因此本章不采用显式瞳孔检测/分割，而是采用轻量端到端建模，让 eye-video 提供 ``是否需要修正'' 与 ``修正方向/幅度'' 的弱监督动态线索。该设计既避免额外标注依赖，也避免三模态在高维特征空间的直接融合造成复杂度与显存膨胀。

本章的三模态输入与对齐方式如下。设 batch size 为 $N$，眼动序列长度 $T_g=100$，world-video 帧数 $T_v=12$，eye-video 帧数 $T_e=42$，图像大小 $H=W=112$。
\begin{equation}
\mathbf{S}\in\mathbb{R}^{N\times T_g\times 2},\qquad \mathbf{S}_{:,t,:}=(x_t,y_t),\ (x_t,y_t)\in[0,1]^2,
\end{equation}
\begin{equation}
\mathbf{V}^w\in\mathbb{R}^{N\times T_v\times 3\times H\times W},\qquad
\mathbf{V}^e\in\mathbb{R}^{N\times T_e\times 1\times H\times W}.
\end{equation}
对齐策略继续以 gaze 时间轴为准：world-video 的对齐沿用上一章 ``凝视引导高斯池化 $\rightarrow$ 时序差分聚合'' 的定义；eye-video 则先在 $T_e$ 上预测补偿序列，再插值到 $T_g$，并用于修正 gaze，形成补偿后的 $\tilde{\mathbf{S}}\in\mathbb{R}^{N\times T_g\times 2}$，再进入与上一章一致的 gaze+world 主干。

\section{多模态融合策略与模型设计}

本章模型记为 \ModelMM{}-Eye，其整体结构可以理解为对上一章 \ModelMM{} 的前置扩展：在原有双模态主干前加入 eye-video 补偿模块，得到补偿 gaze 后，其余模块（gaze 编码、world-video 编码、凝视引导采样、交叉注意力融合、时序建模与分类）均保持与上一章相同的功能路径。为了便于复现，本节将每个模块的输入输出张量维度写清，并对与上一章相同的部分直接引用上一章的公式定义；对新增的 eye-video 补偿分支，则给出完整的结构与公式。

\subsection{Eye-video 补偿分支：从 $\mathbf{S}$ 得到 $\tilde{\mathbf{S}}$（新增部分）}

补偿分支的输入为 eye-video $\mathbf{V}^e\in\mathbb{R}^{N\times T_e\times 1\times 112\times 112}$，输出为对齐到 gaze 时间轴的补偿位移 $\Delta\mathbf{S}\in\mathbb{R}^{N\times T_g\times 2}$ 与门控序列 $\boldsymbol{\gamma}\in\mathbb{R}^{N\times T_g\times 1}$，并生成补偿 gaze $\tilde{\mathbf{S}}\in\mathbb{R}^{N\times T_g\times 2}$。

首先对每帧 eye 图像做轻量卷积编码。实现中卷积通道为 $1\!\rightarrow\!16\!\rightarrow\!32$，步幅为 2，并在末端做全局平均池化得到每帧向量。记该帧级编码器为 $\Phi_e(\cdot)$，则：
\begin{equation}
\mathbf{E}=\Phi_e(\mathbf{V}^e)\in\mathbb{R}^{N\times T_e\times 32},
\end{equation}
其中每个时间步对应一个 32 维特征向量。随后通过线性投影将其映射到补偿分支的内部维度（实现中为 32）：
\begin{equation}
\mathbf{U}_{:,k,:}=\mathbf{E}_{:,k,:}\mathbf{W}_{ep}+\mathbf{b}_{ep},\qquad
\mathbf{U}\in\mathbb{R}^{N\times T_e\times 32}.
\end{equation}

接着使用单向 GRU 在 $T_e$ 维上建模时间动态：
\begin{equation}
\mathbf{H}^e=\mathrm{GRU}_e(\mathbf{U})\in\mathbb{R}^{N\times T_e\times 32}.
\end{equation}
在每个 $k\in\{1,\dots,T_e\}$ 上，分别预测补偿位移（二维）与门控（标量）：
\begin{equation}
\mathbf{o}_k=\mathbf{H}^e_{:,k,:}\mathbf{W}_{\Delta}+\mathbf{b}_{\Delta}\in\mathbb{R}^{N\times 2},\qquad
\mathbf{g}_k=\mathbf{H}^e_{:,k,:}\mathbf{W}_{\gamma}+\mathbf{b}_{\gamma}\in\mathbb{R}^{N\times 1}.
\end{equation}
为了使补偿保持 ``微调'' 属性，对位移采用幅度受限的 $\tanh$ 映射（实现中最大幅度 $m$ 为常数）：
\begin{equation}
\Delta\mathbf{S}^{(e)}_{:,k,:}=m\cdot\tanh(\mathbf{o}_k)\in\mathbb{R}^{N\times 2},\qquad
\boldsymbol{\gamma}^{(e)}_{:,k,:}=\sigma(\mathbf{g}_k)\in\mathbb{R}^{N\times 1}.
\end{equation}
其中 $\sigma(\cdot)$ 为 Sigmoid。由于该输出在 $T_e$ 上定义，而主时间轴为 $T_g$，采用线性插值将序列对齐到 $T_g$：
\begin{equation}
\Delta\mathbf{S}=\mathrm{Interp}_{T_g}\!\left(\Delta\mathbf{S}^{(e)}\right)\in\mathbb{R}^{N\times T_g\times 2},\qquad
\boldsymbol{\gamma}=\mathrm{Interp}_{T_g}\!\left(\boldsymbol{\gamma}^{(e)}\right)\in\mathbb{R}^{N\times T_g\times 1}.
\end{equation}
最后得到补偿 gaze（对每个时间步裁剪到 $[0,1]$）：
\begin{equation}
\tilde{\mathbf{S}}=\mathrm{clip}\left(\mathbf{S}+\boldsymbol{\gamma}\odot\Delta\mathbf{S},\ 0,\ 1\right)\in\mathbb{R}^{N\times T_g\times 2}.
\end{equation}

针对 eye-video 缺失样本，定义可用性标记 $\mathbf{m}\in\{0,1\}^{N\times 1\times 1}$（实现中为 \texttt{eye\_valid}），并将其广播到时间维以旁路补偿：
\begin{equation}
\Delta\mathbf{S}\leftarrow \mathbf{m}\odot\Delta\mathbf{S},\qquad
\boldsymbol{\gamma}\leftarrow \mathbf{m}\odot\boldsymbol{\gamma}.
\end{equation}
当 $\mathbf{m}=0$ 时，有 $\tilde{\mathbf{S}}=\mathbf{S}$，模型自然退化为上一章双模态结构；同时该旁路避免对全零 eye-video 产生无意义的反向缓存，从工程上显著减小显存占用。

\subsection{Gaze 编码器：希尔伯特序编码与动力学注入（沿用上一章公式，补充维度细节）}

补偿后的 gaze 输入为 $\tilde{\mathbf{S}}\in\mathbb{R}^{N\times T_g\times 2}$。其编码流程与上一章 ``希尔伯特序编码器（Hilbert-ordered Gaze Encoder）'' 完全一致：坐标量化、希尔伯特索引构造、哈希映射、索引邻域聚合，以及速度/方向等动力学特征注入均沿用上一章公式定义（包括 $H_t$、$H^{\text{hash}}_t$、邻域集合 $\mathcal{N}_t$、局部嵌入 $\mathbf{e}^{\text{local}}_t$、速度特征 $\mathbf{f}^{\text{vel}}_t=[\Delta x_t,\Delta y_t,v_t,\theta_t]$、坐标嵌入 $\mathbf{e}^{\text{coord}}_t$、以及局部融合 $\mathbf{h}^{\text{local}}_t$ 的构造方式）。本章唯一变化是将上一章中所有 $(x_t,y_t)$ 替换为 $(\tilde{x}_t,\tilde{y}_t)$，即由 $\tilde{\mathbf{S}}$ 生成差分与嵌入。

为明确维度，设特征维度为 $d$（本章实现中 $d=64$），则每个时间步融合后的局部特征为
\begin{equation}
\mathbf{H}^{\text{local}}\in\mathbb{R}^{N\times T_g\times d}.
\end{equation}
上一章在该处使用 TransformerEncoder 做长程建模；本章实现为控制显存，将该序列编码器替换为 1 层双向 GRU（其输入输出长度不变），并加入可学习位置嵌入后线性投影得到 gaze 表征：
\begin{equation}
\mathbf{H}^{\text{gaze}}=\Psi_g(\mathbf{H}^{\text{local}})\in\mathbb{R}^{N\times T_g\times d}.
\end{equation}
其中 $\Psi_g(\cdot)$ 表示 ``位置嵌入 + BiGRU + 线性映射'' 的组合；由于该替换只改变实现层面的序列建模器，不改变上一章的空间编码与动力学注入逻辑，因此本章仍保持 gaze 的处理模式与上一章一致。

\subsection{World-video 编码与凝视引导采样（沿用上一章公式，补充维度细节）}

world-video 输入为 $\mathbf{V}^w\in\mathbb{R}^{N\times T_v\times 3\times 112\times 112}$，输出为与 gaze 时间轴对齐的视频特征 $\mathbf{H}^{\text{video}}\in\mathbb{R}^{N\times T_g\times d}$。该分支由三步组成，均与上一章一致，差异仅在于采样中心改为补偿 gaze $\tilde{\mathbf{S}}$。

\textbf{（1）轻量卷积帧编码器。} 与上一章相同，先对每帧提取特征图。设三层步幅为 2 的卷积编码后特征图大小为 $H_f=W_f=112/8=14$，则：
\begin{equation}
\mathbf{F}=\Phi_w(\mathbf{V}^w)\in\mathbb{R}^{N\times T_v\times d\times H_f\times W_f}
=\mathbb{R}^{N\times T_v\times d\times 14\times 14}.
\end{equation}

\textbf{（2）凝视引导高斯池化。} 该模块的权重定义、归一化方式与 patch 汇聚公式与上一章完全一致（即上一章关于 $w_{s,t_g}(i,j)$、$\hat{w}_{s,t_g}(i,j)$ 与 $\mathbf{p}_{s,t_v,t_g}$ 的三条核心公式）。本章只将权重中心从 $(x_{t_g},y_{t_g})$ 替换为补偿后的 $(\tilde{x}_{t_g},\tilde{y}_{t_g})$。在维度上，若尺度数为 $S$，则 patch 张量为：
\begin{equation}
\mathbf{P}\in\mathbb{R}^{N\times S\times T_v\times T_g\times d}.
\end{equation}
本章实现为控制显存取 $S=1$，因此 $\mathbf{P}\in\mathbb{R}^{N\times 1\times T_v\times T_g\times d}$。

\textbf{（3）时序差分聚合。} 该模块的 ``全局尺度 + 外观/运动建模 + 每尺度融合 + softmax 门控加权'' 的定义与上一章一致（即上一章关于 $\mathbf{g}_{t_v}$、$\mathbf{X}$、$\mathbf{a}_{s,t_g}$、$\mathbf{m}_{s,t_g}$、$\mathbf{u}_{s,t_g}$、$\alpha_{s,t_g}$ 与 $\mathbf{H}^{\text{video}}$ 的公式链）。维度上输出为：
\begin{equation}
\mathbf{H}^{\text{video}}\in\mathbb{R}^{N\times T_g\times d}.
\end{equation}

\subsection{跨模态融合与时序输出（沿用上一章公式，补充维度细节）}

跨模态融合模块与上一章交叉注意力定义一致。输入为
\begin{equation}
\mathbf{H}^{\text{gaze}}\in\mathbb{R}^{N\times T_g\times d},\qquad
\mathbf{H}^{\text{video}}\in\mathbb{R}^{N\times T_g\times d},
\end{equation}
以 gaze 为 Query、video 为 Key/Value，按上一章的多头注意力公式得到 $\mathbf{H}^{\text{attn}}\in\mathbb{R}^{N\times T_g\times d}$，并按上一章残差与层归一化得到融合特征：
\begin{equation}
\mathbf{H}^{\text{fused}}=\mathrm{LayerNorm}\left(\mathbf{H}^{\text{gaze}}+\mathbf{H}^{\text{attn}}\right)\in\mathbb{R}^{N\times T_g\times d}.
\end{equation}

时序建模与分类同样沿用上一章 ``双向 GRU + 线性分类器'' 的形式。设双向 GRU 的隐藏维度为 $d$，则其输出维度为 $2d$：
\begin{equation}
\mathbf{H}^{\text{final}}=\mathrm{BiGRU}(\mathbf{H}^{\text{fused}})\in\mathbb{R}^{N\times T_g\times 2d}.
\end{equation}
再通过全连接映射到类别空间（$C=3$）得到逐时刻 logits：
\begin{equation}
\mathbf{Z}_{:,t,:}=\mathbf{H}^{\text{final}}_{:,t,:}\mathbf{W}_c+\mathbf{b}_c\in\mathbb{R}^{N\times 3},\qquad
\mathbf{W}_c\in\mathbb{R}^{2d\times 3},\ \mathbf{b}_c\in\mathbb{R}^{3}.
\end{equation}
经 softmax 得到类别概率，与上一章定义一致。

\section{实验验证与结果分析}

本章实验沿用上一章的数据划分与训练策略以保证对比公平，核心关注点是：在不改变上一章主干融合机制的前提下，引入 eye-video 运动补偿是否能够稳定提升最终事件检测性能，尤其是对事件边界更敏感的 Event-level 指标是否改善。由于你将补充最终结果与可视化，本节在结构上预留定量表格，同时更充分地给出可视化分析应如何组织、如何解释以及建议呈现的关键证据类型，方便后续直接填图填数。

\subsection{实验设置（与上一章保持一致）}

任务仍为三分类序列标注，类别集合为 GF、GS、GP，样本级与事件级指标的定义与上一章一致。输入片段长度固定为 $T_g=100$，world-video 片段帧数 $T_v=12$，eye-video 帧数 $T_e=42$。训练损失沿用上一章的类别加权交叉熵形式，优化器沿用 AdamW 及相同的学习率与权重衰减设置。为处理 eye-video 缺失，本章在数据加载阶段提供 $\texttt{eye\_valid}$ 标记并在模型内旁路，使缺失样本严格退化为上一章双模态推理路径，从而保证对比不被缺失模态引入的额外噪声干扰。

\subsection{定量结果对比（留白待填）}

表~\ref{tab:mm_eye_final} 用于对比上一章 \ModelMM{}（gaze+world）与本章 \ModelMM{}-Eye（gaze 补偿 + world + eye）。你后续可在表中填入 leave-one-out 汇总结果。

\begin{table*}[htbp]
\centering
\caption{最终结果对比（\%）（留白，后续填入）。上一章为 gaze + world 双模态模型，本章为加入 eye-video 运动补偿的三模态模型。}
\label{tab:mm_eye_final}
\renewcommand{\arraystretch}{1.15}
\setlength{\tabcolsep}{6pt}
\begin{tabular}{lcccccccc}
\toprule
\multirow{2}{*}{方法} & \multicolumn{4}{c}{Sample F1} & \multicolumn{4}{c}{Event F1} \\
\cmidrule(lr){2-5}\cmidrule(lr){6-9}
 & GFi+GFo & GS & GP & W-Avg & GFi+GFo & GS & GP & W-Avg \\
\midrule
上一章：\ModelMM{}（gaze + world） &  &  &  &  &  &  &  &  \\
本章：\ModelMM{}-Eye（gaze 补偿 + world + eye） &  &  &  &  &  &  &  &  \\
\bottomrule
\end{tabular}
\end{table*}

\subsection{结果可视化与定性分析（建议写法与要点）}

为了更有说服力地解释 ``eye-video 为什么有效''，本章建议将可视化分析组织为三类证据，对应补偿分支在系统中的三种作用路径：稳定运动学差分、修正采样锚点、以及改善事件边界一致性。下面给出每类可视化的推荐内容、图的组织方式与应强调的现象，后续你只需替换图与数值即可形成完整分析。

第一类可视化强调 \textbf{补偿对 gaze 轨迹形态与差分特征的影响}。建议在同一段样本上同时绘制原始 gaze $\mathbf{S}$ 与补偿后 gaze $\tilde{\mathbf{S}}$ 的二维轨迹叠加图，并在下方绘制速度曲线 $v_t$（上一章定义的 $v_t=\sqrt{(\Delta x_t)^2+(\Delta y_t)^2}$，这里的差分由 $\mathbf{S}$ 与 $\tilde{\mathbf{S}}$ 分别计算）。需要突出的是：当原始 gaze 存在抖动时，$v_t$ 往往出现密集伪峰值；补偿后 $\tilde{\mathbf{S}}$ 的速度曲线应更 ``稀疏'' 与 ``集中''，尤其在 GS 段体现为更清晰的单峰或少峰结构，在 Fixation 段体现为更低的底噪。对应图可以包括：(1) 轨迹图；(2) $v_t$ 对比；(3) $\|\Delta\mathbf{S}_t\|$ 与 $\gamma_t$ 的时间曲线，用于说明补偿不是全程强干预，而是在特定时段激活。

第二类可视化强调 \textbf{补偿对 world-video 采样锚点与证据提取的影响}。上一章的凝视引导高斯池化以 gaze 坐标为中心形成权重热图并在特征图上加权汇聚。本章建议在若干关键帧上，将采样热图叠加到原始 world-video 帧或特征图可视化上，分别展示以 $\mathbf{S}$ 与 $\tilde{\mathbf{S}}$ 为中心的热图差异，并配合说明 ``错误锚点会采样到无关区域''。在场景包含局部运动目标或显著目标时，这种差异最直观：补偿后采样中心更贴近目标区域，视频分支提取到的运动证据更一致，从而使后续交叉注意力更容易在正确位置上建立匹配。对应图可以组织为三列：原始帧、原始 gaze 热图、补偿 gaze 热图；并选择包含明显运动或遮挡的片段作为示例。

第三类可视化强调 \textbf{事件级边界与预测序列的一致性}。建议将同一段样本的真实标签序列与模型预测序列对齐展示（例如用颜色条或逐点类别曲线），并同时给出上一章模型与本章模型的预测对比。分析时应重点观察：GS 的起止边界是否更贴近标注，GP 段是否更连贯，是否减少 ``碎片化'' 的短事件插入。由于 Event F1 对边界与片段一致性更敏感，这类可视化可以直接解释定量提升来源。若你还会展示典型失败案例，也建议在失败案例中对比 $\gamma_t$ 的激活情况：如果在困难片段 $\gamma_t$ 接近 0，说明补偿分支可能未能提供有效修正；如果 $\gamma_t$ 过高但仍失败，则可能提示 eye-video 质量不佳或场景变化导致主干误判，这有助于后续改进方向的讨论。

在写作层面，建议将可视化分析明确与模型结构对应起来：先用一段话说明 ``本章的改动只在补偿前端，主干与上一章一致，因此定性差异主要来自 $\tilde{\mathbf{S}}$ 的变化''；再按上述三类图逐一说明现象与原因，最后回到结果表强调补偿带来的整体收益。这样能够形成从结构--现象--指标的闭环，避免可视化停留在展示层面。

\section{本章小结}

本章在上一章 gaze+world 双模态融合框架基础上，引入 eye-video 构建三模态眼动事件检测方法 \ModelMM{}-Eye。引入 eye-video 的核心理由在于：双模态框架对场景证据建模充分，但对 gaze 观测误差缺乏直接修复通道，而 gaze 既承担运动学判别信号又承担视频采样锚点，抖动与漂移会同时破坏两条路径。为避免三模态直接高维融合带来的复杂度膨胀，本章将 eye-video 的作用限定为前置运动补偿：通过轻量卷积编码与 GRU 时序建模预测幅度受限的补偿位移与门控系数，得到补偿后的 gaze $\tilde{\mathbf{S}}$，并将其输入到与上一章一致的 gaze 编码、world-video 凝视引导采样、交叉注意力融合与时序分类器中。实验部分将通过最终定量对比与可视化分析共同验证该补偿策略在边界一致性与整体识别鲁棒性上的增益来源。


\chapter{总结与展望}

\section{研究总结}

本研究围绕眼动事件检测的核心问题，从单模态时频融合、空间-时序融合到多模态自适应融合，构建了完整的技术体系，主要完成以下工作：

1. 双流时频眼动事件检测：针对全局依赖捕捉与局部边界定位的矛盾，提出\ModelTF{}框架，通过Mamba网络建模长时全局时序依赖，CNN提取局部时频特征，注意力机制实现双流融合。实验表明，\ModelTF{}在三个公开数据集上均取得最优性能，事件级Kappa系数最高达0.975，接近人类专家水平，且计算复杂度低，适合实时应用。

2. 眼动序列空间编码研究：针对现有方法忽略空间结构特征的问题，提出三种空间编码方案（时间窗口矩阵、注视点热力图、时空图结构），设计STFN空间-时序融合模型。实验验证，时空图结构编码性能最优，STFN模型相比基准模型，PSO检测精度提升2.9\%，丰富了眼动事件的特征表示维度。

3. 多模态眼动事件检测：针对复杂场景下的抗干扰需求，构建多模态数据集（眼动坐标、瞳孔直径、场景图像、传感器数据），提出MMEMD自适应多模态融合框架。实验表明，MMEMD模型的平均F1值达到0.962，PSO检测精度提升6.3\%，且具有更强的抗干扰鲁棒性，为实际应用提供了可靠解决方案。

4. 技术体系构建与验证：从单模态到多模态，从时序-时频特征到空间-时序-语义特征，构建了层次化、模块化的眼动事件检测技术体系，通过大量实验验证了各模块的有效性和兼容性，为不同应用场景提供灵活选择。

\section{研究创新点总结}

1. 提出“全局时序-局部时频”双流融合框架\ModelTF{}，首次将Mamba网络应用于眼动信号建模，实现线性时间复杂度的长时依赖捕捉，解决了传统模型的效率与性能矛盾；

2. 设计时空图结构编码方法，将眼动序列建模为时空图，通过GNN提取拓扑结构特征，实现空间-时序特征的深度融合，提升了对复杂轨迹模式的识别能力；

3. 提出自适应多模态融合机制，通过模态注意力动态调整权重，跨模态交互层捕捉非线性关联，解决了模态异质性和数据缺失问题，显著提升了模型的鲁棒性；

4. 构建了完整的眼动事件检测技术体系，涵盖单模态、双模态、多模态等多种形态，为学术研究和工程应用提供了全面的技术支撑。

\section{研究局限与未来展望}

研究局限
1. PSO事件检测精度仍有提升空间：PSO信号微弱、持续时间短，现有模型对其多模态特征的捕捉仍不充分；
2. 空间编码的动态适应性不足：当前空间编码方案基于固定参数（如窗口大小、网格尺寸），难以适应不同个体、不同任务的眼动轨迹差异；
3. 多模态数据的完整性和标注质量有待提升：现有多模态数据集的场景多样性和噪声类型不够丰富，影响模型的泛化能力；
4. 模型的实时性与部署优化不足：多模态模型的计算复杂度相对较高，需要进一步优化以适应移动设备部署需求。

未来展望
1. 精细化PSO事件检测：结合生物力学知识，设计专用的多尺度、多模态PSO特征提取模块，提升微弱信号的识别能力；
2. 自适应空间编码方法：引入元学习、强化学习等技术，动态调整空间编码参数，适应不同个体和任务的差异；
3. 大规模多模态数据集构建：收集更多真实场景下的多模态眼动数据（如驾驶、医疗诊断场景），丰富噪声类型和场景多样性；
4. 轻量化模型设计与部署：采用模型压缩、量化、剪枝等技术，优化多模态模型的计算复杂度，实现移动设备的实时部署；
5. 跨领域迁移学习研究：探索眼动事件检测模型在不同领域（如神经诊断、人机交互、自动驾驶）的迁移学习方法，降低标注成本；
6. 端到端多任务学习：将眼动事件检测与特征提取、意图识别、疾病诊断等任务联合训练，构建一体化眼动分析系统。

\end{document}